\documentclass[12pt, a4paper]{article} 

\usepackage[utf8]{inputenc}       
\usepackage[T1]{fontenc}          
\usepackage[brazil]{babel}        
\usepackage{geometry}          
\usepackage{graphicx}             
\usepackage{array}              
\usepackage{hyperref}             
\usepackage{amsmath}
\usepackage{lastpage}
\usepackage[export]{adjustbox} 
\usepackage{fancyhdr}
\usepackage{tikz}
\usetikzlibrary{shapes, arrows.meta, positioning, fit, shapes.geometric, backgrounds}
\usepackage{tabularx}
\usepackage{array}
\usepackage{multirow}
\usepackage{appendix} 
\usepackage{longtable} 
\usepackage{url}
\usepackage{listings}
\usepackage{xcolor}
\usepackage[table]{colortbl}
\usepackage{tcolorbox}
\usepackage{float}

\pagestyle{fancy}           
\fancyhf{}                  
\renewcommand{\headrulewidth}{0pt} 
\renewcommand{\footrulewidth}{0.4pt} 

\fancyfoot[L]{\small \quad Vers. : \docVersion} 
\fancyfoot[R]{\small Pagina : \thepage{} / \pageref{LastPage}} 

\fancypagestyle{plain}{
  \fancyhf{}
  \renewcommand{\headrulewidth}{0pt}
  \renewcommand{\footrulewidth}{0.4pt}
  \fancyfoot[L]{\small \quad Vers. : \docVersion}
  \fancyfoot[R]{\small Pagina : \thepage{} / \pageref{LastPage}}
}

\geometry{a4paper, top=2.5cm, bottom=2.5cm, left=2.5cm, right=2.5cm}

% === INFORMACOES DO DOCUMENTO ===
\newcommand{\projectTitle}{CAOS Agent}
\newcommand{\projectSubtitle}{Framework Agêntico \& Governança}
\newcommand{\docType}{Documento de Especificação Técnica} 
\newcommand{\docVersion}{18.2}
\newcommand{\docDate}{\today} 
\newcommand{\docCreationDate}{23 de janeiro de 2026} 
\newcommand{\companyUnit}{VivAIOT - Ecossistema de Monitoramento IoT} 
\newcommand{\companyAddress}{Sistema de Gestão de Ativos}
\newcommand{\addressedTo}{Equipe de Desenvolvimento Viva AI} 
\newcommand{\authorName}{Lucas Severino}

% === CORES ===
\definecolor{vivacolor}{RGB}{0,102,204}
\definecolor{caoscolor}{RGB}{128,0,128} % Roxo para o CAOS
\definecolor{codegreen}{rgb}{0,0.6,0}
\definecolor{codegray}{rgb}{0.5,0.5,0.5}
\definecolor{codepurple}{rgb}{0.58,0,0.82}
\definecolor{backcolour}{rgb}{0.95,0.95,0.92}

% === CONFIGURACAO DE CODIGO ===
\lstdefinestyle{mystyle}{
    backgroundcolor=\color{backcolour},   
    commentstyle=\color{codegreen},
    keywordstyle=\color{magenta},
    numberstyle=\tiny\color{codegray},
    stringstyle=\color{codepurple},
    basicstyle=\ttfamily\footnotesize,
    breakatwhitespace=false,         
    breaklines=true,                 
    captionpos=b,                    
    keepspaces=true,                 
    numbers=left,                    
    numbersep=5pt,                  
    showspaces=false,                
    showstringspaces=false,
    showtabs=false,                  
    tabsize=2
}
\lstset{
    style=mystyle,
    inputencoding=utf8,
    extendedchars=true,
    literate={á}{{\'a}}1 {ã}{{\~a}}1 {é}{{\'e}}1 {ç}{{\c{c}}}1 {ó}{{\'o}}1 {õ}{{\~o}}1 {í}{{\'i}}1 {ú}{{\'u}}1 {Á}{{\'A}}1 {Ã}{{\~A}}1 {É}{{\'E}}1 {Ç}{{\c{C}}}1 {Ó}{{\'O}}1 {Õ}{{\~O}}1 {Í}{{\'I}}1 {Ú}{{\'U}}1 {â}{{\^a}}1 {ê}{{\^e}}1 {ô}{{\^o}}1 {Â}{{\^A}}1 {Ê}{{\^E}}1 {Ô}{{\^O}}1
}

% === CAIXAS ===
\newtcolorbox{notabox}[1][]{
    colback=caoscolor!5!white,
    colframe=caoscolor,
    title=#1,
    fonttitle=\bfseries
}

\newtcolorbox{avisobox}[1][]{
    colback=yellow!10!white,
    colframe=orange!80!black,
    title=#1,
    fonttitle=\bfseries
}

\begin{document}

% ============================================================================
% CAPA
% ============================================================================
\begin{titlepage}
    \centering 

    \noindent 
    \begin{tabular}[b]{@{}l@{}}
        \includegraphics[height=3cm]{img/logo_vivaiot.png}
    \end{tabular}
    \hfill
    \begin{tabular}[b]{@{}r@{}} 
        \small 
        Vers. : \docVersion \\
        Data : \docDate \\
        Pagina : \thepage{} / \pageref{LastPage} \\
        Data de criação : \docCreationDate \\
    \end{tabular}

    \vspace{2cm} 

    \large{\textbf{VivAIOT - Sistema de Monitoramento de Ativos IoT}} \\ 

    \normalsize{\companyUnit} \\ 

    \vspace{3cm} 

    {\Huge\bfseries \docType} \\ 

    \vspace{1cm} 

    {\Large\bfseries \projectTitle} \\ 
    {\large \projectSubtitle} \\ 
    \vspace{2cm} 

    \begin{flushleft} 
        Destinado a : \\
        \addressedTo \\ 
    \end{flushleft}

    \vfill 

    \noindent 
    \fbox{ 
        \begin{minipage}{0.95\textwidth}
            \begin{minipage}[t]{0.5\textwidth}
                \raggedright
                \textbf{Redigido por :} \\
                \vspace{0.5cm} 
                \authorName \\
            \end{minipage}
            \begin{minipage}[t]{0.5\textwidth}
                \raggedleft
            \end{minipage}
            \vspace{1cm} 
        \end{minipage}
    }

    \vspace{1cm}

    \centering
    \small

    \companyAddress \\ 

\end{titlepage}

% ============================================================================
% SUMARIO
% ============================================================================
\newpage
\tableofcontents
\newpage

\listoffigures 
\newpage 

% ============================================================================
% HISTORICO DE REVISOES
% ============================================================================
\section*{Histórico de Revisões}
\vspace{0.5cm} 

\begin{center} 
    \begin{tabular}{|m{1.5cm}|m{3cm}|m{7.5cm}|m{2.5cm}|} 
        \hline
        \textbf{Versão} & \textbf{Data} & \textbf{Modificações} & \textbf{Autor} \\
        \hline
        1.0 & 23/01/2026 & Versão inicial do documento & \authorName \\
        \hline
        5.0 & 24/01/2026 & Adição da arquitetura LangGraph e Temporal.io & \authorName \\
        \hline
        10.0 & 25/01/2026 & Definição do algoritmo de veredito e faixas de decisão & \authorName \\
        \hline
        14.0 & 27/01/2026 & Motor de Guardrails e políticas JSON & \authorName \\
        \hline
        15.0 & 28/01/2026 & Catálogo completo de 38 Guardrails em 6 categorias & \authorName \\
        \hline
        16.0 & 28/01/2026 & Fluxos, contratos, observabilidade, testes e recovery & \authorName \\
        \hline
        17.0 & 28/01/2026 & Revisão estrutural: seções, tools expandidas & \authorName \\
        \hline
        17.5 & 28/01/2026 & Correções arquiteturais: 6 etapas, Oracle topology, Calibração, Extensibilidade & \authorName \\
        \hline
        18.0 & 28/01/2026 & Glossário expandido, Oracle no diagrama, seção de Limitações & \authorName \\
        \hline
        18.1 & 28/01/2026 & Diagrama de sequência com todos os 4 agentes + CAOS & \authorName \\
        \hline
        18.2 & 28/01/2026 & CARE como executor, 6 fases visuais no diagrama & \authorName \\
        \hline
        18.3 & 01/02/2026 & Substituição do Temporal.io por Google Cloud Workflows & \authorName \\
        \hline
    \end{tabular}
\end{center} 
\newpage 

% ============================================================================
% GLOSSARIO
% ============================================================================
\section*{Glossário} 
\begin{description} 
    \item[ASM :] Agentic State Machine - Máquina de Estados Agêntica
    \item[CAOS :] Centralized Autonomous Operating System - Sistema Operacional Autônomo Centralizado
    \item[Digital Twin :] Representação virtual de um ativo físico, sincronizada em tempo real
    \item[FSM :] Finite State Machine - Máquina de Estados Finita
    \item[Guardrail :] Regra de segurança que valida ações antes da execução
    \item[HAL :] Hardware Abstraction Layer - Camada que traduz comandos lógicos em protocolos industriais
    \item[Human-in-the-Loop :] Processo que exige aprovação humana
    \item[LangGraph :] Framework para construir agentes com grafos e estado
    \item[LOTO :] Lockout/Tagout - Protocolo de segurança que bloqueia equipamentos durante manutenção
    \item[mTLS :] Mutual TLS - Autenticação mútua via certificados entre cliente e servidor
    \item[Oracle :] Agente de predição que projeta cenários futuros usando modelos de ML
    \item[Orquestração :] Coordenação centralizada de agentes e serviços
    \item[Google Pub/Sub :] Sistema de mensageria para comunicação assíncrona global e serverless
    \item[RAG :] Retrieval-Augmented Generation - Técnica que combina busca vetorial com geração de texto
    \item[SLA :] Service Level Agreement - Acordo de Nível de Serviço
    \item[Viva Co-Pilot :] Interface humana para interação e aprovação
    \item[WORM :] Write Once Read Many - Armazenamento imutável para auditoria legal
\end{description}
\newpage

% ============================================================================
% SECAO 1: CONTEXTO E DEFINICAO CORE
% ============================================================================
% ============================================================================
% CAPITULO 1: FUNDAÇÃO E METODOLOGIA
% ============================================================================
\section{Contexto e Objetivos}
\label{sec:fundacao}

O \textbf{CAOS (Centralized Autonomous Operating System)} representa uma mudança paradigmática na automação industrial: a transição de sistemas reativos (determinísticos) para sistemas cognitivos (probabilísticos).

\subsection{O Conceito Judge-LM}

Diferente de controladores tradicionais (PIDs, PLCs) que operam em lógica linear, o CAOS atua como um \textbf{Juiz Sintético}. Sua função não é apenas executar scripts, mas sim \textit{arbitrar} conflitos entre múltiplas fontes de dados (Sensores, Manuais, Previsões e Regras de Negócio) para determinar a ação ótima em cenários de alta incerteza.

\subsection{Topologia do Sistema}

A arquitetura do sistema é baseada em \textbf{Loops de Controle Cognitivo}. Conforme ilustrado na Figura \ref{fig:arquitetura_macro}, o sistema é composto por quatro domínios funcionais integrados via barramento de baixa latência:

\begin{figure}[h!]
    \centering
    \resizebox{\textwidth}{!}{%
    \begin{tikzpicture}[
        node distance=1.0cm and 0.5cm, % Reduced Horizontal spacing
        auto,
        block/.style={
            rectangle, 
            draw=blue!60, 
            fill=blue!5, 
            thick, 
            text width=2.5cm, % Slightly narrower blocks
            align=center, 
            rounded corners, 
            minimum height=1.5cm
        },
        cloud/.style={
            ellipse, 
            draw=gray!40, 
            fill=gray!5, 
            dashed, 
            minimum height=1.5cm,
            text width=2.2cm,
            align=center
        },
        line/.style={
            draw, 
            -Latex, 
            thick
        }
    ]

    % Nodes
    \node [cloud] (world) {Mundo Físico\\(Fábrica)};
    
    % Sensors and Manuals stacked
    \node [block, right=of world, yshift=1.2cm] (sensors) {Sensores \&\\Telemetria};
    \node [block, right=of world, yshift=-1.2cm] (manuals) {Manuais \&\\Contexto};
    
    % Core centrally placed to the right of inputs
    \node [block, right=of sensors, yshift=-1.2cm, minimum height=3.5cm, text width=3.2cm, fill=blue!10] (core) {\textbf{CAOS CORE}\\(Judge-LM)};
    
    \node [block, right=of core] (hal) {HAL\\(Care)};
    
    \node [cloud, right=of hal] (actuators) {Atuadores\\(Motores)};

    % Labels inside Core
    \node [font=\scriptsize, below=0.2cm of core.north] {Decision Engine};

    % Paths
    \path [line] (world.north) |- (sensors.west);
    \path [line] (world.south) |- (manuals.west);
    
    \path [line] (sensors) -- node[above, font=\scriptsize, sloped] {Dados} (core.160);
    \path [line] (manuals) -- node[below, font=\scriptsize, sloped] {Metadados} (core.200);
    
    \path [line] (core) -- node[midway, font=\scriptsize, align=center] {Comando} (hal);
    \path [line] (hal) -- node[midway, font=\scriptsize, align=center] {Protocolo} (actuators);
    
    % Feedback loop
    \draw [line, dashed] (actuators.south) |- ++(0,-1.0) -| (world.south) node[pos=0.25, below, font=\scriptsize] {Feedback};

    \end{tikzpicture}
    }
    \caption{Diagrama de Blocos: Loop Cibernético de Controle}
    \label{fig:arquitetura_macro}
\end{figure}

\begin{enumerate}
    \item \textbf{Contexto (ATLAS)}: Ingestão de estado contextual. Fornece ao CAOS o \textit{significado} dos valores (ex: especificações técnicas indicam que 90°C é crítico para \textit{este} modelo específico).
    \item \textbf{Instinto (SENTINEL)}: Detecção de anomalias em tempo real (<10ms). Detecta \textit{quando} os limites foram violados e classifica a severidade do alerta.
    \item \textbf{Predição (ORACLE)}: Projeção de futuros possíveis. Agente \textbf{Opcional/Condicional} para otimização de custos de tokens. Calcula probabilidades: "Se executar ação X, qual a chance de resultado Y?".
    \item \textbf{Deliberação (CAOS Core)}: O núcleo de inferência. Consome dados de ATLAS, SENTINEL e (se necessário) ORACLE para calcular o Veredito matemático ($V$).
    \item \textbf{Verificação (GUARDRAILS)}: Sistema Imunológico. Veta ou aprova a ação proposta com base em políticas de segurança.
    \item \textbf{Atuação (CARE)}: Camada de atuação ("os músculos") motorizada pelo \textbf{Google Cloud Workflows}. Traduz decisões lógicas em protocolos industriais seguros via HAL.
\end{enumerate}

\subsection{A Lógica de Julgamento (The Judgment Loop)}

O CAOS opera em um loop contínuo de percepção e decisão. Diferente de sistemas baseados em regras rígidas ("Se/Então"), ele calcula um \textbf{Veredito Matemático} ($V$) para cada frame de realidade:

\begin{equation}
V = \frac{\text{Benesses (Atlas + } \hat{\text{O}}\text{racle)} - \text{Riscos (Sentinel)}}{\text{Guardrails (Leis)}}
\end{equation}

\noindent \textit{Nota: Esta é uma representação conceitual. Para a equação multi-dimensional exata e calibração de pesos, consulte a Seção \ref{sec:verdict_formula}.}

O processo ocorre em \textbf{6 etapas cíclicas}:

\begin{table}[h!]
\centering
\small
\begin{tabular}{|c|l|l|l|}
\hline
\textbf{Etapa} & \textbf{Agente} & \textbf{Função} & \textbf{Domínio} \\ \hline
1. Contexto & ATLAS & Coleta fatos (Digital Twin, manuais, specs) & Memória \\ \hline
2. Instinto & SENTINEL & Detecta anomalias e alertas & Percepção \\ \hline
3. Simulação & ORACLE & Projeta futuros possíveis (\textbf{Opcional/ROI}) & Predição \\ \hline
4. Julgamento & CAOS & Aplica fórmula do Veredito ($V$) & Decisão \\ \hline
5. Verificação & GUARDRAILS & Veta ou aprova a ação proposta & Segurança \\ \hline
6. Execução & CARE & Executa ação física (HAL → Device) & Atuação \\ \hline
\end{tabular}
\caption{Ciclo Cognitivo de 6 Etapas}
\end{table}

\subsection{Topologia de Fluxo de Dados (Arquitetura Híbrida)}

O CAOS utiliza uma arquitetura híbrida onde os \textbf{Guardrails funcionam como uma constituição interna} — um conjunto de princípios que o CAOS consulta para decidir qual ação tomar. Isso significa que, se a IA não alucinar, a ação proposta já foi verificada internamente e está ``dentro da lei''.

Para ações de \textbf{alto risco}, o sistema adiciona camadas de verificação externa:
\begin{itemize}
    \item \textbf{Risco Baixo (LOW)}: Ação vai diretamente para o CARE.
    \item \textbf{Risco Médio (MEDIUM)}: Verificação automática adicional (regras de negócio, limites).
    \item \textbf{Risco Alto (HIGH)}: Requer aprovação humana explícita (Human-in-the-Loop).
\end{itemize}

\begin{figure}[h!]
    \centering
    \resizebox{\textwidth}{!}{%
    \begin{tikzpicture}[
        node distance=1.2cm and 1.5cm,
        % Styles
        agent/.style={rectangle, draw=black, thick, fill=pink!30, minimum width=2.2cm, minimum height=0.9cm, align=center, font=\footnotesize},
        core/.style={rectangle, draw=purple!80, thick, fill=purple!20, minimum width=3cm, minimum height=3cm, align=center, font=\footnotesize\bfseries},
        guardrail/.style={rectangle, draw=purple!60, thick, fill=purple!10, minimum width=2.2cm, minimum height=0.7cm, align=center, font=\scriptsize},
        decision/.style={diamond, draw=orange!80, thick, fill=orange!10, aspect=1.8, font=\scriptsize\bfseries, inner sep=1pt},
        verify/.style={rectangle, draw=blue!80, thick, fill=blue!10, minimum width=1.8cm, minimum height=0.7cm, align=center, font=\scriptsize},
        hitl/.style={rectangle, draw=orange!80, thick, fill=orange!20, minimum width=1.8cm, minimum height=0.7cm, align=center, font=\scriptsize},
        exec/.style={rectangle, draw=green!80, thick, fill=green!20, minimum width=2.2cm, minimum height=1.2cm, align=center, font=\footnotesize\bfseries},
        device/.style={ellipse, draw=gray!60, dashed, fill=gray!5, minimum width=1.8cm, minimum height=0.9cm, align=center, font=\footnotesize},
        arrow/.style={->, thick, >=stealth},
        dashedarrow/.style={->, thick, >=stealth, dashed}
    ]
    
    % === INPUT AGENTS (LEFT) ===
    \node [agent] (atlas) at (0, 1.5) {ATLAS\\(Contexto)};
    \node [agent] (sentinel) at (0, 0) {SENTINEL\\(Instinto)};
    \node [agent] (oracle) at (0, -1.5) {ORACLE\\(Predição)};
    
    % === CAOS CORE WITH GUARDRAILS INSIDE ===
    \node [core] (caos) at (4, 0) {};
    \node [font=\footnotesize\bfseries, purple!80] at (4, 0.8) {CAOS};
    \node [font=\scriptsize] at (4, 0.3) {(Judge)};
    \node [guardrail] (guardrails) at (4, -0.6) {GUARDRAILS\\(Constituição)};
    
    % === RISK CLASSIFICATION ===
    \node [decision] (risk) at (7.5, 0) {Risco?};
    
    % === VERIFICATION PATHS ===
    \node [verify] (autoverify) at (10, 0) {Verificação\\Auto};
    \node [hitl] (human) at (10, -2) {HITL\\(Humano)};
    
    % === EXECUTION ===
    \node [exec] (care) at (13, 0) {CARE\\(Execução)};
    \node [device] (device) at (16, 0) {Device};
    
    % === ARROWS: Agents -> CAOS ===
    \draw [arrow] (atlas.east) -- node[above, font=\tiny] {$A$} (caos.155);
    \draw [arrow] (sentinel.east) -- node[above, font=\tiny] {$S$} (caos.180);
    \draw [dashedarrow] (oracle.east) -- node[above, font=\tiny] {$O$} (caos.205);
    
    % === ARROWS: CAOS -> Risk Decision ===
    \draw [arrow] (caos.east) -- node[above, font=\tiny] {Ação Candidata} (risk.west);
    
    % === ARROWS: Risk Paths ===
    % LOW - Direct to CARE
    \draw [arrow, green!60!black] (risk.north) -- ++(0, 1.2) -| node[pos=0.15, above, font=\tiny, green!60!black] {LOW} (care.north);
    
    % MEDIUM - Auto Verification
    \draw [arrow, blue!60!black] (risk.east) -- node[above, font=\tiny, blue!60!black] {MEDIUM} (autoverify.west);
    \draw [arrow] (autoverify.east) -- (care.west);
    
    % HIGH - Human in the Loop
    \draw [arrow, orange!80!black] (risk.south) -- node[left, font=\tiny, orange!80!black] {HIGH} (human.north);
    \draw [arrow] (human.east) -| (care.south);
    
    % === ARROWS: CARE -> Device ===
    \draw [arrow] (care.east) -- node[above, font=\tiny] {Command} (device.west);
    
    % === FEEDBACK LOOP ===
    \draw [dashedarrow, gray] (device.south) -- ++(0,-1.5) -| node[below, font=\tiny, pos=0.15] {Feedback} (caos.south);
    
    \end{tikzpicture}
    }
    \caption{Topologia de Fluxo Híbrido: Agentes → CAOS (c/ Guardrails) → Classificação de Risco → Execução}
    \label{fig:data_flow}
\end{figure}

\noindent (Para a definição matemática exata e pesos, consulte o \textbf{Capítulo \ref{sec:judge_algo}}).

\subsection{Stack Tecnológica de Suporte}

Apesar da mudança conceitual, a infraestrutura base permanece robusta para suportar a "Consciência":

% ============================================================================
% CAPITULO 2: ARQUITETURA MACRO E STACK
% ============================================================================
\section{Arquitetura}
\label{sec:arquitetura_macro}



\subsection{Decisões Tecnológicas (Rationale)}
\begin{itemize}
    \item \textbf{Cérebro (LangGraph on Cloud Run)}: O LangGraph é o \textit{Framework de Lógica} que gerencia o estado e loops cognitivos. Ele roda empacotado dentro de containers serverless no \textbf{Google Cloud Run}, garantindo escalabilidade.
    \item \textbf{Inteligência (Vertex AI)}: O LangGraph não "pensa", ele "orquestra". O pensamento real vem do \textbf{Gemini 1.5 Pro}, acessado via API do Vertex AI.
    \item \textbf{Memória (Firestore)}: O estado dos agentes é persistido no Firestore, permitindo que o Cloud Run escale a zero sem perder o contexto (Memory Saver).
    \item \textbf{Atuação (Agente CARE via Cloud Workflows)}: O \textbf{Google Cloud Workflows} é o motor que dá vida ao Agente CARE. Ele recebe as decisões do CAOS e as transforma em uma sequência resiliente de comandos para os drivers do HAL, garantindo retentativas e tratamento de erros sem sobrecarregar a lógica cognitiva do CAOS.
\end{itemize}

\subsection{Infraestrutura Atual}
Abaixo, o fluxo de dados entre os componentes da stack:

\begin{figure}[h!]
    \centering
    \resizebox{\textwidth}{!}{%
    \begin{tikzpicture}[
        node distance=2.0cm,
        auto,
        >=stealth,
        % Styles
        container/.style={
            rectangle, draw=black!50, dashed, fill=white, inner sep=0.3cm, rounded corners,
            label={[anchor=north west, font=\tiny, color=gray]north west:#1}
        },
        component/.style={
            rectangle, draw=black, thick, fill=white, 
            minimum width=2.5cm, minimum height=1.0cm, align=center, font=\footnotesize
        },
        db/.style={
            cylinder, shape border rotate=90, draw=black, thick, aspect=0.25,
            minimum width=1.5cm, minimum height=1.5cm, align=center, font=\footnotesize
        },
        bus/.style={
            rectangle, draw=red!80, fill=red!5, thick,
            minimum width=12cm, minimum height=0.6cm, align=center, font=\bfseries\small
        }
    ]

    % BUS (Pub/Sub)
    \node [bus] (bus) {GOOGLE PUB/SUB (Global Event Bus)};

    % BRAIN CLUSTER (LangGraph)
    \node [component, above=2.0cm of bus.center, fill=blue!10] (cortex) {\textbf{Cortex}\\(LangGraph)};
    \node [component, left=0.8cm of cortex, fill=blue!5] (sense) {\textbf{Sense}\\(Fusion)};
    \node [component, right=0.8cm of cortex, fill=blue!5] (act) {\textbf{Act}\\(Planning)};
    
    % Container drawn on background layer to avoid covering nodes
    \begin{pgfonlayer}{background}
        \node [container={LangGraph Runtime}, fit=(sense)(cortex)(act), fill=blue!2] (brain) {};
    \end{pgfonlayer}

    % MEMORY (Postgres)
    \node [db, above=1.5cm of cortex, fill=gray!10] (pg) {\textbf{Postgres}\\(Vector Store\\+ Checkpointer)};

    % MOTOR (CARE Agent / Cloud Workflows)
    \node [component, below=2.5cm of bus.center, fill=orange!10, text width=4.5cm] (care_wf) {\textbf{CARE Agent}\\(Cloud Workflows)};
    \node [component, right=1.5cm of care_wf, fill=green!10] (hal_exec) {\textbf{HAL Drivers}\\(Cloud Run)};

    % IOT SENSOR NODE - positioned to the left of CARE to avoid overlap
    \node [component, left=1.5cm of care_wf, fill=gray!10] (iot_sensor) {\textbf{Sensor IoT}\\(Edge)};

    % CONNECTIONS
    % Telemetry arrow - from IoT Sensor to Bus (left side, no overlap)
    \draw [->, thick, red] (iot_sensor.north) -- node[midway, left, font=\tiny, text=red] {Telemetry (100Hz)} (iot_sensor.north |- bus.south);
    
    \draw [->, thick] (bus.north -| sense.south) -- node[left, font=\tiny, pos=0.5] {Pull} (sense.south);
    \draw [->, thick] (act.south) -- node[right, font=\tiny, pos=0.5] {Push Event} (bus.north -| act.south);
    
    \draw [<->, dashed] (cortex) -- node[right, font=\tiny] {RAG / State} (pg);

    % 4. Bus -> CARE -> HAL
    \draw [->, thick] (bus.south) -- node[right, font=\tiny] {Trigger CARE} (care_wf.north);
    \draw [->, thick] (care_wf) -- node[above, font=\tiny] {HTTP/LOTO} (hal_exec);

    % 5. HAL -> Physical Effect -> IoT Sensor (feedback loop)
    \draw [->, thick, dashed] (hal_exec.south) -- ++(0,-0.8) -| (iot_sensor.south) node[pos=0.25, below, font=\tiny] {Effect (Side-Effect)};

    \end{tikzpicture}
    }
    \caption{Arquitetura de Sistemas: Cérebro (CAOS) e Músculos (CARE)}
    \label{fig:stack_flow}
\end{figure}

A tabela abaixo detalha cada escolha:

\begin{table}[h!]
\centering
\renewcommand{\tabularxcolumn}[1]{m{#1}} % Force vertical centering
\begin{tabularx}{\textwidth}{|>{\centering\arraybackslash}m{1.5cm}|l|l|X|}
\hline
\textbf{Logo} & \textbf{Categoria} & \textbf{Tecnologia} & \textbf{Função} \\ \hline
\includegraphics[height=0.8cm]{img/langgraph-color.png} & Orquestração & LangGraph & Loop de Inferência Cognitiva (O Cérebro) \\ \hline
\includegraphics[height=0.8cm]{img/vertexai.png} & Mensageria & \textbf{Google Cloud Pub/Sub} & Sistema Nervoso (Ingestão de Eventos) \\ \hline
\includegraphics[height=0.8cm]{img/vertexai.png} & Instinto & \textbf{Google Pub/Sub} & Canal de Ingestão de Alertas (via Sentinel) \\ \hline
\includegraphics[height=0.8cm]{img/vertexai.png} & Inteligência & \textbf{Vertex AI} & API de Inferência (Cérebro Auxiliar) \\ \hline
\includegraphics[height=0.8cm]{img/vertexai.png} & Atuação & \textbf{CARE (via Workflows)} & Execução Serverless e Orquestração de HAL \\ \hline
\end{tabularx}
\caption{Infraestrutura Cognitiva}
\end{table}

\subsection{Análise de Redundância: LangGraph vs Cloud Workflows}
Uma dúvida comum na arquitetura do CAOS é a aparente redundância entre usar dois orquestradores. Ambos gerenciam grafo e estado, mas operam em domínios de execução distintos.

\subsubsection{Separação de Responsabilidades (Cérebro vs Operação)}

\begin{table}[h!]
\centering
\begin{tabular}{|l|p{6cm}|p{6cm}|}
\hline
\textbf{Característica} & \textbf{LangGraph (Cérebro)} & \textbf{Cloud Workflows (Músculo)} \\ \hline
\textbf{Foco Principal} & Raciocínio, Contexto e Decisão. & Orquestração Serverless e Chamadas de API. \\ \hline
\textbf{Tipo de Estado} & \textbf{Cognitivo (Short-Lived)}. Mantém o histórico da conversa e do pensamento atual. & \textbf{Operacional (Event-Driven)}. Gerencia o fluxo de execução de passos. \\ \hline
\textbf{Loop de Controle} & Dinâmico/Probabilístico (LLM decide o próximo passo). & Sequencial/Determinístico (YAML define os gatilhos e passos). \\ \hline
\textbf{Exemplo de Uso} & "Devo desligar este compressor com base nestes 3 manuais?" & "Recebi ordem de desligar. Enviando sinal HTTP e verificando feedback em 5s." \\ \hline
\end{tabular}
\caption{Matriz de Responsabilidade: LangGraph vs Cloud Workflows}
\end{table}

\subsubsection{Resolução de Conflitos}
Para evitar "dois chefes" mandando na mesma tarefa, estabelecemos a regra de ouro:
\begin{itemize}
    \item \textbf{LangGraph é o Legislador}: Ele define \textbf{O QUE} deve ser feito. Ele não executa a ação, ele emite a "Ordem de Serviço".
    \item \textbf{Cloud Workflows é o Executor}: Ele define \textbf{COMO} e \textbf{QUANDO} a ação será concluída. Ele recebe a ordem e coordena as chamadas técnicas aos drivers do CARE.
\end{itemize}

Esta separação resolve o problema de concorrência/redundância ao desacoplar estritamente a decisão da execução.
% ============================================================================
% 2.3 Detalhamento da Stack
% ============================================================================
\subsection{Detalhamento da Implementação por Tecnologia}

Esta seção destrincha o papel e a implementação de cada componente da stack:

\subsubsection{LangGraph on Cloud Run (O Cérebro)}
\textbf{Implementação}: Utiliza o conceito de \textit{StateGraph} onde cada nó é uma função Python pura que recebe e retorna um estado tipado (`JudgeState`).
\textbf{Uso}: Centraliza a lógica de decisão. Diferente de uma "Chain" linear, o Grafo permite ciclos, rodando em containers \textit{stateless} (Cloud Run) que escalam conforme a demanda cognitiva.
\begin{lstlisting}[language=Python, basicstyle=\tiny\ttfamily]
workflow = StateGraph(JudgeState)
workflow.add_node("sense", sense_node)
workflow.add_node("oracle", oracle_node)
workflow.add_node("cortex", cortex_node)
workflow.add_node("guardrails", guardrails_node)
workflow.add_node("act", act_node)
workflow.set_entry_point("sense")
\end{lstlisting}

\subsubsection{Firestore (Interface de Memória)}
\textbf{Papel}: O CAOS não gerencia o banco de dados diretamente. Ele consome o estado do ativo (Digital Twin) mantido pelo agente Atlas via leituras no Firestore. O CAOS trata o Firestore como uma API de "Leitura de Estado".

\subsubsection{Google Pub/Sub (O Sistema Nervoso)}
\textbf{Implementação}: Utiliza tópicos globais dessacoplados. Sensores enviam JSON para `projects/vivaiot/topics/telemetry`.
\textbf{Uso}: Gatilho primário. Garante que se o Cérebro estiver "dormindo" (Cold Start), a mensagem fica na fila (persistência de 7 dias) e nenhum sinal vital é perdido.

\subsubsection{Firestore (A Memória)}
\textbf{Implementação}:
\begin{itemize}
    \item \textbf{Gêmeo Digital}: Banco NoSQL Realtime. Cada ativo é um documento: `assets/pump-01`.
    \item \textbf{Eventos}: O Firestore dispara um evento para o Cloud Run sempre que um campo crítico (ex: `temp`) muda.
\end{itemize}
\textbf{Uso}: "Qual o estado atual?". O Agente não pergunta ao sensor, ele pergunta ao Firestore (Single Source of Truth).

\subsubsection{Sentinel Subscription (Interface de Instinto)}
\textbf{Integração}: O CAOS é agnóstico à implementação do Sentinel. Ele apenas assina o tópico \texttt{sentinel.alerts} e espera um JSON validado. O processamento de stream ocorre fora do perímetro do CAOS.

\subsubsection{Google Vertex AI (O Córtex)}
\textbf{Uso}: Onde o raciocínio acontece. O LangGraph estrutura os dados e invoca o Vertex AI para obter o "Veredito". O CAOS utiliza o modelo Gemini 1.5 Pro configurado para determinismo (Temperature=0).

\newpage

% ============================================================================
% CAPITULO 3: MAQUINA DE ESTADOS (FSM)
% ============================================================================
\section{Máquina de Estados}
\label{sec:fsm}

O CAOS opera sob uma Máquina de Estados Finita (FSM) global que governa o ciclo de vida de cada ticket ou anomalia:

\begin{enumerate}
    \item \textbf{IDLE}: Estado de repouso (Zero custo de computação).
    \item \textbf{TRIGGERED}: Recebimento de evento via Sentinel ou Oracle.
    \item \textbf{PERCEIVING}: Data Fusion (Atlas + Contexto).
    \item \textbf{PREDICTING}: Simulação via Oracle (Opcional).
    \item \textbf{JUDGING}: Raciocínio Híbrido (Analyze Impacts $\rightarrow$ Calculate S \& V $\rightarrow$ Check Guardrails $\rightarrow$ Classify Risk).
    \item \textbf{EXECUTING}: Módulo Care envia comando ao hardware.
    \item \textbf{RESOLVED}: Sucesso confirmado.
    \item \textbf{BLOCKED}: Guardrail impediu ação.
    \item \textbf{ESCALATED}: Humano convocado.
\end{enumerate}

\begin{figure}[h!]
    \centering
    \begin{tikzpicture}[
        node distance=2cm, 
        auto, 
        thick, 
        scale=0.8, 
        every node/.style={transform shape},
        state/.style={circle, draw=blue!50, fill=blue!5, thick, minimum size=1.2cm, font=\scriptsize}
    ]
        \node[state] (idle) {IDLE};
        \node[state] (trig) [right of=idle] {TRIGGER};
        \node[state] (perceive) [right of=trig] {PERCEIVE};
        \node[state] (judge) [below of=perceive] {JUDGE};
        \node[state] (exec) [left of=judge] {EXECUTE};
        
        \draw[->] (idle) -- (trig);
        \draw[->] (trig) -- (perceive);
        \draw[->] (perceive) -- (judge);
        \draw[->] (judge) -- (exec);
        \draw[->] (exec) -| (idle);
    \end{tikzpicture}
    \caption{Ciclo de Vida Global do CAOS}
    \label{fig:global_fsm}
\end{figure}

\subsection{O Grafo Cognitivo (Agentic State Machine - ASM)}
\label{sec:cognitive_graph}

A inteligência do CAOS não reside em um único prompt longo, mas em um \textbf{Estado Agêntico Orquestrado} via LangGraph. O sistema opera como um grafo cíclico dirigido onde o estado flui e evolui entre nós especializados.

\subsubsection{1. O Contrato de Dados (BrainState)}
O Grafo compartilha um único objeto de estado imutável que persiste no Firestore:
\begin{itemize}
    \item \texttt{input\_event}: O gatilho vindo do Sentinel.
    \item \texttt{context\_chunks}: Memórias recuperadas pelo Atlas.
    \item \texttt{reasoning\_trace}: A trilha de pensamento acumulada.
    \item \texttt{safety\_score}: Validação contínua dos Guardrails.
\end{itemize}

\subsubsection{2. Definição dos Nós (Estações de Pensamento)}
Cada nó representa uma competência cognitiva específica:

\begin{enumerate}
    \item \textbf{Nó: Sense (Percepção)}: Responsável pelo \textit{Data Fusion}. Consolida alertas do Sentinel com o contexto do Atlas (Manuais, Contratos e Digital Twin).
    \item \textbf{Nó: Oracle (Simulação)}: Nó opcional (bypassável em \textit{Fast-Track}). Realiza projeções de impacto financeiro e técnico.
    \item \textbf{Nó: Cortex (Julgamento)}: O "Core" cognitivo. Executa o raciocínio híbrido em 4 etapas: (1) Análise de Dimensões, (2) Cálculo de $S$ e $V$, (3) Aplicação de Guardrails contextuais e (4) Classificação de Risco. Propõe a ação final.
    \item \textbf{Nó: Guardrails (Enforcement)}: Nó determinístico de segurança final. Valida a ação proposta contra as leis imutáveis. Em caso de violação, força o \textbf{Veto} ou o \textit{Reciclo} (re-planejamento).
    \item \textbf{Nó: Act (Despacho)}: Finaliza o estado da decisão e publica o comando de execução para o Agente CARE (Cloud Workflows).
\end{enumerate}

\subsubsection{3. Fluxo e Bordas Condicionais}
O grafo utiliza \textbf{Conditional Edges} para decidir o caminho dinamicamente, priorizando a segurança sobre a latência ou o custo:
\begin{itemize}
    \item \texttt{Sense} $\rightarrow$ \texttt{Cortex} (\textbf{Bypass Oracle}): Ativado em \textbf{Risco Crítico}. O sistema prioriza a latência zero para proteção imediata do ativo, ignorando simulações futuras.
    \item \texttt{Sense} $\rightarrow$ \texttt{Oracle} (\textbf{Otimização}): Ativado em \textbf{Risco Baixo/Médio}. O sistema utiliza o tempo disponível para projetar cenários e otimizar custos operacionais.
    \item \texttt{Guardrails} $\rightarrow$ \texttt{Cortex} (\textbf{Reciclo}): Se a proposta for vetada, o fluxo retorna para re-planejamento.
\end{itemize}

\begin{figure}[h!]
    \centering
    \resizebox{\textwidth}{!}{%
    \begin{tikzpicture}[
        node distance=1.5cm,
        block/.style={rectangle, draw=blue!50, thick, fill=blue!5, minimum width=2.5cm, minimum height=0.8cm, align=center, rounded corners},
        decision/.style={diamond, draw=red!50, thick, fill=red!5, aspect=2, font=\tiny\bfseries},
        arrow/.style={-Latex, thick}
    ]
        \node [block] (sense) {Sense};
        \node [decision, right=of sense] (risk) {Crítico?};
        \node [block, below=1.2cm of risk] (oracle) {Oracle};
        \node [block, right=of risk] (cortex) {Cortex};
        \node [block, right=of cortex] (guard) {Guardrails};
        \node [block, right=of guard] (act) {Act};

        \draw [arrow] (sense) -- (risk);
        \draw [arrow] (risk) -- node[above, font=\tiny] {Sim} (cortex);
        \draw [arrow] (risk) -- node[left, font=\tiny] {Não} (oracle);
        \draw [arrow] (oracle) -| (cortex);
        \draw [arrow] (cortex) -- (guard);
        \draw [arrow] (guard.north) -- ++(0,0.5) -| node[pos=0.25, above, font=\tiny] {Veto (Reciclo)} (cortex.north);
        \draw [arrow] (guard) -- node[above, font=\tiny] {Ok} (act);
    \end{tikzpicture}
    }
    \caption{Fluxo de Orquestração Cognitiva (LangGraph)}
    \label{fig:langgraph_flow}
\end{figure}
O CAOS não é um monólito. Ele é composto por serviços desacoplados que escalam independentemente:

\begin{enumerate}
    \item \textbf{Ingestion Service (Sentinel)}: O "pulmão" do sistema. Recebe milhares de mensagens de telemetria por segundo, filtra ruídos e descarta dados inúteis antes de acionar o cérebro.
    \item \textbf{Cognitive Service (Cortex)}: Onde roda o LangGraph. Este serviço escala baseado em consumo de CPU/Tokens. É "Stateless" (o estado fica no Postgres).
    \item \textbf{Action Service (Care)}: Google Cloud Workflows. Responsável pelos efeitos colaterais. Se a rede cair, o sistema retenta a orquestração com backoff exponencial.
\end{enumerate}

Essa separação garante que um pico de tráfego de IoT (Sentinel) não derrube a capacidade de raciocínio (Cortex) nem a execução de tarefas críticas (Care).

% ============================================================================
% CAPITULO 3: O CEREBRO (ORQUESTRACAO COGNITIVA)
% ============================================================================
\section{Orquestrador (LangGraph)}
\label{sec:brain}

O fluxo de pensamento do Judge-LM é modelado como um grafo direcionado cíclico (LangGraph). A Figura \ref{fig:multi_agent} ilustra o papel do CAOS como o maestro desta orquestra cognitiva:

\begin{figure}[h!]
    \centering
    \resizebox{0.9\textwidth}{!}{%
    \begin{tikzpicture}[
        node distance=2.5cm,
        auto,
        font=\small\sffamily,
        % Styles
        hub/.style={
            circle, draw=blue!80, fill=blue!10, very thick, 
            minimum size=3.5cm, align=center, text width=2.5cm, font=\bfseries
        },
        agent/.style={
            rectangle, draw=black!60, fill=white, thick, 
            minimum width=3cm, minimum height=1.5cm, align=center, rounded corners
        },
        arrow/.style={->, >=stealth, thick, color=gray!80, text=black}
    ]

    % CENTRAL HUB
    \node [hub] (caos) {CAOS\\Orchestrator\\(Judge-LM)};

    % SATELLITE AGENTS
    \node [agent, above left=of caos, fill=purple!5] (atlas) {\textbf{ATLAS}\\(Memória \& Contexto)};
    \node [agent, above right=of caos, fill=red!5] (sentinel) {\textbf{SENTINEL}\\(Instinto \& Alerta)};
    \node [agent, below left=of caos, fill=cyan!5] (oracle) {\textbf{ORACLE}\\(Predição Futura)};
    \node [agent, below right=of caos, fill=orange!5] (care) {\textbf{CARE}\\(Ação \& HAL)};

    % CONNECTIONS (Star Topology)
    
    % Atlas <-> CAOS
    \draw [arrow] (caos.135) -- node[midway, right, font=\tiny] {Query} (atlas.315);
    \draw [arrow] (atlas.300) -- node[midway, left, font=\tiny] {RAG} (caos.150);
    
    % Sentinel <-> CAOS
    \draw [arrow] (sentinel.240) -- node[midway, right, font=\tiny] {Alert} (caos.30);
    \draw [arrow] (caos.45) -- node[midway, left, font=\tiny] {Ack} (sentinel.225);

    % Oracle <-> CAOS
    \draw [arrow, dashed] (caos.225) -- node[midway, left, font=\tiny] {Simulate (On-demand)} (oracle.45);
    \draw [arrow, dashed] (oracle.60) -- node[midway, right, font=\tiny] {Forecast} (caos.210);

    % Care <-> CAOS
    \draw [arrow] (caos.315) -- node[midway, left, font=\tiny] {Plan} (care.135);
    \draw [arrow] (care.120) -- node[midway, right, font=\tiny] {Result} (caos.330);

    \end{tikzpicture}
    }
    \caption{Topologia Estrela: Orquestração Multi-Agente}
    \label{fig:multi_agent}
\end{figure}

\subsection{Contratos de Dados e Tipos}

Para garantir a interoperabilidade entre os 5 agentes, o CAOS utiliza schemas Pydantic estritos.

\subsubsection{1. Entrada (TriggerPayload)}
O dado que "acorda" o CAOS via Pub/Sub. Todo agente (Sentinel, Oracle, Co-Pilot) deve respeitar este contrato:
\begin{lstlisting}[language=Python]
class TriggerPayload(BaseModel):
    event_id: str = Field(..., description="UUID único do evento")
    source: Literal["sentinel", "oracle", "manual"]
    timestamp: datetime
    
    # [POR QUE DICT?] Flexibilidade para diferentes fontes.
    # Ex Sentinel: {"sensor": "chiller-01", "temp": 95.5}
    # Ex Oracle: {"forecast_metric": "failure_prob", "value": 0.89}
    payload: Dict[str, Any] 
    
    # Metadados de rastreabilidade (Tenant, Asset ID)
    context: Dict[str, str]
\end{lstlisting}

\subsubsection{2. Memória de Trabalho (JudgeState)}
O estado compartilhado no LangGraph, enriquecido passo a passo:
\begin{lstlisting}[language=Python]
class JudgeState(TypedDict):
    # Identidade e Gatilho
    trigger: TriggerPayload
    
    # Camada 1 - Sense (Percepção)
    atlas_context: Optional[AssetContext]        # Digital Twin + Manuais + Contratos
    risk_assessment: Optional[RiskAssessment]    # Sentinel Data
    
    # Camada 2 - Cortex (Raciocínio & Decisão)
    oracle_forecast: Optional[Prediction]        # Simulação Oracle (se houver)
    
    # [NOVO] Dimensões e Escores (Fórmula 4.5)
    risk_dimensions: Dict[str, float]            # {R_F, R_Fin, R_C, R_K}
    severity_score: float                        # Escore S
    verdict_score: float                         # Escore V (-1.0 a +1.0)
    
    # Explicação e Ação
    reasoning_trace: List[str]                   # Chain-of-Thought
    proposed_action: Optional[ActionSchema]      # Ação gerada pelo LLM
    guardrail_violations: List[str]              # Lista de IDs (ex: PHYS_001)
    is_fast_track: bool                          # Se Oracle foi bypassado
\end{lstlisting}

\subsubsection{3. Saída (ActionSchema)}
O comando final enviado ao Google Cloud Workflows para execução física:
\begin{lstlisting}[language=Python]
class ActionSchema(BaseModel):
    action_type: Literal["shutdown", "notification", "ticket"]
    target_resource: str  # ex: "chiller-04"
    
    # [POR QUE DICT?] Cada ação exige parâmetros únicos.
    # Ex Shutdown: {"force": True, "delay_seconds": 0}
    # Ex Ticket: {"priority": "High", "assignee": "maintenance_team"}
    parameters: Dict[str, Any] 
    
    # Justificativa dada pelo LLM para auditoria
    reasoning: str
    confidence: float
\end{lstlisting}

\subsection{Nós do Pensamento}
\begin{enumerate}
    \item \textbf{sensory\_integration}: Coleta dados do Atlas e Sentinel (Data Fusion).
    \item \textbf{frontal\_cortex}: Nó \textbf{condicional}. Invoca o Oracle apenas se a severidade for menor que CRITICAL ou se houver alta incerteza no veredito prévio, economizando tokens em ações de segurança imediata.
    \item \textbf{conscious\_judgement}: Aplica a Equação do Veredito.
    \item \textbf{motor\_planning}: Converte decisão em sinais para o HAL.
\end{enumerate}

\subsection{Data Fusion (Percepção Unificada)}
O nó \textit{sensory\_integration} realiza a fusão vetorial necessária para o Oracle e para o Veredito:

\begin{equation}
S_{t} = \Phi(Sentinel_{raw}) \oplus \Psi(Atlas_{context})
\end{equation}

E a projeção oracular que depende deste estado:

\begin{equation}
O_{forecast} = \Omega(S_t, \text{Proposed Action})
\end{equation}

Onde cada operador possui uma função epistemológica específica:

\begin{itemize}
    \item \textbf{$S_t$} (State at $t$): O estado unificado da realidade no tempo $t$.
    \item \textbf{$\Phi$ (Normalização)}: Função que padroniza inputs brutos (ex: converte $mV$ para $^\circ C$).
    \item \textbf{$\Psi$ (Contextualização)}: Função que injeta metadados estáticos (ex: limites de garantia).
    \item \textbf{$\Omega$ (Projeção)}: Função oracular que estima $S_{t+1}$ baseada em tendências.
    \item \textbf{$\oplus$ (Merge Seguro)}: Operador de união imutável com validação de schema Pydantic.
\end{itemize}

Se o Sentinel reportar "Vibração Alta" ($\Phi$) mas o Atlas informar "Manutenção Agendada" ($\Psi$), a fusão ($\oplus$) resulta logicamente em \textit{Maintenance State}, desativando alertas de pânico.

\subsection{Exemplo de Cadeia de Raciocínio (Chain-of-Thought)}
Abaixo, um log real extraído do \textit{frontal\_cortex} durante um incidente de alta temperatura:

\begin{lstlisting}[language=Python, caption=Trace de Raciocínio do Judge-LM, basicstyle=\tiny\ttfamily]
{
  "thought_process": [
    {
      "step": 1,
      "observation": "Input Sensor: Temp=92C on Chiller-04",
      "memory_recall": "Atlas says MaxOpTemp=90C. Warranty void if >95C."
    },
    {
      "step": 2,
      "hypothesis": "If I do nothing, asset likely fails in <5min.",
      "risk_assessment": "Sentinel reports P1 (Critical). People nearby."
    },
    {
      "step": 3,
      "tool_check": "Can I use 'emergency_shutdown'? Yes. Guardrails allow P1.",
      "alternative": "Can I use 'fan_boost'? No, fan is already at 100%."
    },
    {
      "step": 4,
      "verdict": {
        "action": "EMERGENCY_SHUTDOWN",
        "confidence": 0.98,
        "justification": "Imminent hardware destruction. Cost of downtime < Cost of asset."
      }
    }
  ]
}
\end{lstlisting}

\newpage

% ============================================================================
% CAPITULO 4: O JUIZ (ALGORITMO DE DECISAO)
% ============================================================================
\subsection{O Juiz: Algoritmo de Decisão}
\label{sec:judge_algo}

A decisão final ($V$) não é fruto de aleatoriedade, mas sim de uma \textbf{Balança Matemática Rigorosa}.

\begin{figure}[h!]
    \centering
    \begin{tikzpicture}[
        node distance=1.5cm,
        auto,
        startstop/.style={rectangle, rounded corners, minimum width=3cm, minimum height=1cm, text centered, draw=black, fill=red!10},
        process/.style={rectangle, minimum width=3cm, minimum height=1cm, text centered, draw=black, fill=orange!10},
        decision/.style={diamond, aspect=2, minimum width=3cm, minimum height=1cm, text centered, draw=black, fill=green!10},
        arrow/.style={thick,->,>=stealth}
    ]

    \node (inputs) [startstop] {Inputs (A, S)};
    \node (decide_oracle) [decision, below=of inputs] {Exige Oracle?};
    \node (oracle_sim) [process, right=of decide_oracle, xshift=0.5cm] {ORACLE ($O$)};
    \node (weights) [process, below=of decide_oracle] {Pesos ($W_x$)};
    \node (sum) [process, below=of weights] {Soma ($V_{raw}$)};
    \node (guard) [decision, below=of sum] {Guardrails?};
    
    % VETO to the right
    \node (veto) [startstop, right=of guard, xshift=0.5cm] {VETO};
    
    \node (classify) [process, below=of guard] {Classificação};
    \node (action) [startstop, below=of classify] {Ação Final};

    \draw [arrow] (inputs) -- (decide_oracle);
    \draw [arrow] (decide_oracle) -- node[above] {Sim} (oracle_sim);
    \draw [arrow] (decide_oracle) -- node[right] {Não ($O=1$)} (weights);
    \draw [arrow] (oracle_sim) |- (weights);
    \draw [arrow] (weights) -- (sum);
    \draw [arrow] (sum) -- (guard);
    
    \draw [arrow] (guard) -- node[above] {Não} (veto);
    \draw [arrow] (guard) -- node[right] {Sim} (classify);
    \draw [arrow] (classify) -- (action);
    
    \end{tikzpicture}
    \caption{Fluxograma Lógico do Algoritmo de Decisão}
    \label{fig:veredicto}
\end{figure}

Conforme a Figura \ref{fig:veredicto}, o processo de julgamento coloca em pratos opostos os benefícios da ação e os riscos detectados.
No prato esquerdo, acumulamos os \textbf{Fatos} (dados do Atlas) e \textbf{Benefícios} (previsões de lucro do Oracle). No prato direito, empilhamos os \textbf{Riscos} (alertas do Sentinel).
O resultado dessa pesagem deve, obrigatoriamente, passar pelo funil dos \textbf{Guardrails}. Se qualquer regra de segurança for maior que o benefício, o funil bloqueia a saída, resultando em um VETO inquestionável.

A equação formal, que expande a visão simplificada apresentada no Capítulo 1, é dada por:

\begin{equation}
V = \frac{(W_A \cdot Norm(A)) + (W_O \cdot \hat{O}) - (W_S \cdot (1 + S))}{1 + \sum (P_G \cdot V_G)}
\end{equation}

Onde o termo oracular $\hat{O}$ é definido de forma binária conforme a necessidade de simulação:

\begin{equation}
\hat{O} = 
\begin{cases} 
Norm(O) & \text{se } S < \text{THRESHOLD}_{critical} \text{ (Prioridade: Inteligência/Otimização)} \\
1.0 & \text{se } S \geq \text{THRESHOLD}_{critical} \text{ (Prioridade: Sobrevivência/Latência)}
\end{cases}
\end{equation}

Onde cada componente pondera a decisão final em um espaço vetorial:

\begin{itemize}
    \item \textbf{$V$} (Veredito Final): Score normalizado entre $-1.0$ (Veto) e $+1.0$ (Ação Prioritária).
    \item \textbf{$W_A, W_O, W_S$} (Pesos): Configurações ajustáveis de importância (ex: $W_S=0.3$ equilibra segurança).
    \item \textbf{$Norm(A), \hat{O}$}: Funções de normalização que convertem a saída qualitativa do LLM (baseada em Atlas) e a simulação do Oracle em um escalar $[0, 1]$.
    \item \textbf{$S$} (Severidade): Grau de risco medido pelo Sentinel (0=Seguro, 1=Catastrófico).
    \item \textbf{$P_G \cdot V_G$} (Penalidade de Guardrail): Se um guardrail ativo ($V_G=1$) for violado com alta penalidade ($P_G$), o denominador cresce tendendo ao infinito, zerando $V$.
\end{itemize}

\subsection{Calculadora de Cenários (Exemplos Reais)}
Para validar a equação, submetemos 5 cenários de borda ao sistema. Considere $W_A=0.6, W_O=0.6, W_S=0.3$.

\begin{table}[h!]
\centering
\begin{tabular}{|l|c|c|c|c|c|l|}
\hline
\textbf{Cenário} & \textbf{A} & \textbf{O} & \textbf{S} & \textbf{Guard} & \textbf{$V_{calc}$} & \textbf{Ação} \\ \hline
1. Normal Operation & 1.0 & 0.8 & 0.0 & OK (1.0) & \textbf{0.78} & \textit{Executar} \\ \hline
2. Fogo (Compressor) & 0.2 & 0.0 & \textbf{1.0} & OK (1.0) & \textbf{-0.48} & \textbf{SHUTDOWN} \\ \hline
3. Lucro Alto/Risco Médio & 1.0 & 0.9 & 0.4 & OK (1.0) & \textbf{0.72} & \textit{Executar} \\ \hline
4. Bloqueio Físico (Porta) & 1.0 & 1.0 & 0.0 & \textbf{VETO} & \textbf{ERR} & \textbf{BLOCK} \\ \hline
5. Override Humano & 1.0 & 1.0 & 0.8 & OK (1.0) & \textbf{0.66} & \textit{Sugerir} \\ \hline
\end{tabular}
\caption{Matriz de Validação da Equação de Veredito}
\label{tab:calc_cenarios}
\end{table}

\noindent \textbf{Análise do Cenário 2 (Fogo)}:
Mesmo sem benefício (O=0), a alta severidade do Sentinel ($S=1.0$) empurra o score para negativo ($0.12 - 0.6 = -0.48$), o que cai na faixa de \textbf{VETO/SHUTDOWN} imediato, ignorando qualquer tentativa de "otimização".

\subsection{Definição das Faixas de Decisão}
Com base nos cálculos acima, calibramos os seguintes thresholds definitivos para o Judge-LM:

\begin{table}[h!]
\centering
\begin{tabular}{|c|l|l|}
\hline
\textbf{Score ($V$)} & \textbf{Decisão (Enum)} & \textbf{Comportamento do Sistema} \\ \hline
$V \leq 0.0$ & \textbf{BLOCKED} & Bloqueio imediato. Nenhuma ação é tomada. \\ \hline
$0.0 < V \leq 0.3$ & \textbf{ALERT} & Baixa confiança. Notifica operador sem sugerir ação. \\ \hline
$0.3 < V \leq 0.7$ & \textbf{SUGGEST} & Confiança moderada. Sugere ação para aprovação humana. \\ \hline
$V > 0.7$ & \textbf{EXECUTE} & Alta confiança. Executa ação autonomamente (se $P_G=0$). \\ \hline
\end{tabular}
\caption{Tabela de Faixas de Decisão (Decision Bands)}
\label{tab:decision_thresholds}
\end{table}

\subsection{Validação e Calibração de Pesos}

Os pesos $W_A$, $W_O$ e $W_S$ não são arbitrários. Eles são calibrados através de um processo rigoroso:

\begin{enumerate}
    \item \textbf{Backtesting Histórico}: Aplicar a fórmula em incidentes passados e verificar se o Veredito $V$ teria gerado a ação correta. Métricas: precisão, recall, F1-score.
    \item \textbf{A/B Testing em Staging}: Rodar duas versões do sistema com pesos diferentes e medir outcomes (tempo de resolução, custo de downtime, satisfação do cliente).
    \item \textbf{RLHF (Reinforcement Learning from Human Feedback)}: Quando operadores sobrescrevem decisões do CAOS, o sistema registra o delta e ajusta os pesos via gradiente.
    \item \textbf{Análise de Sensibilidade}: Variar cada peso individualmente ($\pm 0.1$) e observar o impacto na distribuição de decisões.
\end{enumerate}

\begin{notabox}[Pesos Iniciais (Heurísticos)]
$W_A = 0.6$, $W_O = 0.6$, $W_S = 0.3$ — Equilíbrio entre autonomia e segurança. Estes valores são ponto de partida e devem ser calibrados por vertical (varejo vs indústria) e por tenant.
\end{notabox}

\subsection{Integração com Vertex AI (Onde o LLM Pensa)}

A fórmula $V$ é calculada por código determinístico, mas os \textbf{scores de entrada} ($A$, $O$) e a \textbf{proposição da ação} são gerados pelo LLM (Gemini via Vertex AI). A tabela abaixo detalha os pontos exatos de chamada:

\begin{table}[h!]
\centering
\small
\begin{tabular}{|c|l|l|}
\hline
\textbf{Etapa} & \textbf{Executor} & \textbf{O que acontece} \\ \hline
1. Alerta & SENTINEL (Código) & Detecta anomalia, envia $S$ \\ \hline
2. Contexto & ATLAS (RAG) & Retorna manuais e specs \\ \hline
3. Propor Ação & \textbf{Vertex AI (LLM)} & "Recomendo SHUTDOWN" \\ \hline
4. Calcular $A$ & \textbf{Vertex AI (LLM)} & "Manual apoia? A=0.9" \\ \hline
6. Simular & \textbf{ORACLE (Opcional)} & Recebe $S_t$ se necessário para ROI/Simulação \\ \hline
6. Calcular $V$ & CÓDIGO (Python) & $V = 0.6(0.9) + 0.6(0.7) - 0.3(1.0)$ \\ \hline
7. Guardrails & CÓDIGO (Regras) & Verifica violações \\ \hline
8. Explicar & \textbf{Vertex AI (LLM)} & Gera justificativa textual \\ \hline
\end{tabular}
\caption{Pontos de Integração com Vertex AI}
\end{table}

A Figura \ref{fig:llm_flow} ilustra a posição do LLM no fluxo de decisão:

\begin{figure}[h!]
    \centering
    \begin{tikzpicture}[
        node distance=1.5cm and 2cm,
        box/.style={rectangle, draw=black, thick, minimum width=2.2cm, minimum height=1cm, align=center, font=\footnotesize, rounded corners=2pt},
        llm/.style={rectangle, draw=blue!80, thick, fill=blue!10, minimum width=3cm, minimum height=1.2cm, align=center, font=\footnotesize\bfseries, rounded corners=3pt},
        formula/.style={rectangle, draw=orange!80, thick, fill=orange!10, minimum width=5cm, minimum height=1cm, align=center, font=\footnotesize, rounded corners=3pt},
        guard/.style={rectangle, draw=red!60, thick, fill=red!5, minimum width=2.5cm, minimum height=0.8cm, align=center, font=\footnotesize, rounded corners=2pt},
        arrow/.style={->, thick, >=stealth},
        curved/.style={->, thick, >=stealth, bend left=20},
        curvedR/.style={->, thick, >=stealth, bend right=20}
    ]
    
    % Row 1: Inputs (spread out more)
    \node [box, fill=gray!5] (sentinel) {SENTINEL\\$S$};
    \node [box, fill=gray!5, right=2.5cm of sentinel] (atlas) {ATLAS\\Contexto};
    \node [box, fill=gray!5, right=2.5cm of atlas] (oracle) {ORACLE\\Previsão};
    
    % Row 2: LLM (centered under Atlas)
    \node [llm, below=2cm of atlas] (llm) {VERTEX AI\\(Gemini LLM)\\Propor Ação, Calcular $A$};
    
    % Row 3: Formula (centered, wider)
    \node [formula, below=1.8cm of llm] (formula) {FÓRMULA: $V = W_A \cdot A + W_O \cdot O - W_S \cdot (1+S)$};
    
    % Row 4: Guardrails
    \node [guard, below=1.5cm of formula] (guard) {GUARDRAILS};
    
    % Arrows with curves
    % Oracle -> LLM (curved left)
    \node [left=0.5cm of llm, font=\tiny, text=blue] (fused) {Fused $(S, A)$};
    \draw [arrow, blue] (sentinel.south) to[bend right] (llm.north west);
    \draw [arrow, blue] (atlas.south) -- (llm.north);
    \draw [arrow, red, thick, dashed] (llm.east) to[bend left] node[right, font=\tiny] {Simulate (Opcional)} (oracle.south);
    \draw [arrow, forestgreen, thick, dashed] (oracle.west) to[bend right] node[above, font=\tiny] {Forecast} (llm.south east);
    
    % Oracle -> Formula (dashed, O goes direct)
    \draw [curved, dashed, gray] (oracle.south) to[out=-90, in=30] node[right, font=\tiny, pos=0.7] {$O$} (formula.east);
    
    % LLM -> Formula (straight, with A, Ação)
    \draw [arrow] (llm.south) -- node[right, font=\tiny] {$A$, Ação} (formula.north);
    
    % Sentinel -> Formula (curved left, bypass LLM - S goes direct)
    \draw [curvedR, dashed, gray] (sentinel.south) to[out=-90, in=150] node[left, font=\tiny, pos=0.7] {$S$} (formula.west);
    
    % Formula -> Guardrails
    \draw [arrow] (formula.south) -- node[right, font=\tiny] {$V$} (guard.north);
    
    \end{tikzpicture}
    \caption{Posição do LLM no Fluxo de Decisão}
    \label{fig:llm_flow}
\end{figure}

\begin{avisobox}[Importante]
O LLM \textbf{não calcula} a fórmula $V$. Ele \textbf{propõe a ação} e \textbf{avalia conformidade} com documentos. A fórmula serve como validação matemática auditável do raciocínio do LLM.
\end{avisobox}


\subsection{Thresholds Configuráveis por Contexto}

As faixas de decisão \textbf{não são universais}. Elas podem ser customizadas por múltiplas dimensões:

\begin{table}[h!]
\centering
\small
\begin{tabular}{|l|l|l|}
\hline
\textbf{Dimensão} & \textbf{Exemplo} & \textbf{Impacto} \\ \hline
Tenant & Cliente conservador & EXECUTE só se $V > 0.85$ \\ \hline
Tipo de Ativo & Freezer de sorvete vs Chiller & Thresholds térmicos diferentes \\ \hline
Modo de Operação & Horário comercial vs noturno & Tolerância maior à noite \\ \hline
Contrato & SLA Premium vs Basic & Velocidade de resposta diferente \\ \hline
Região & Brasil vs Europa & Regulamentações diferentes \\ \hline
\end{tabular}
\caption{Dimensões de Customização de Thresholds}
\end{table}

A configuração é feita via JSON no nível do Tenant:

\begin{lstlisting}[language=Python, caption=Exemplo de Configuração de Thresholds, basicstyle=\tiny\ttfamily]
{
  "tenant_id": "cliente_xyz",
  "decision_bands": {
    "BLOCKED": { "max": 0.0 },
    "ALERT": { "min": 0.0, "max": 0.25 },
    "SUGGEST": { "min": 0.25, "max": 0.85 },
    "EXECUTE": { "min": 0.85 }
  },
  "weights": { "W_A": 0.25, "W_O": 0.25, "W_S": 0.50 }
}
\end{lstlisting}

\subsection{Dinâmica de Julgamento Não-Linear (O Fim do If/Else)}
A principal inovação do CAOS é a transição da \textbf{Lógica Booleana (Clássica)} para o \textbf{Controle Cognitivo (Fuzzy)}.

\paragraph{O Problema da Lógica Rígida}
Em sistemas tradicionais, o mundo é binário. Um sensor acusando $95.1^\circ C$ dispara um alarme se o limite for $95.0^\circ C$. Toda a complexidade da realidade (nuances, contextos, exceções) precisa ser hardcoded em milhares de linhas de `if/else`, tornando a manutenção impossível.

\paragraph{A Solução: Espaço Vetorial de Decisão}
O CAOS não "testa condições". Ele "pesa princípios".
O Veredito $V$ não é um interruptor, é um gradiente.
\begin{enumerate}
    \item \textbf{Recuperação Semântica (Atlas)}: Em vez de consultar variáveis, o agente lê manuais e normas. Ele entende que "96ºC é aceitável por 10 min apenas durante o degelo". Isso transforma \textit{Dados} em \textit{Conhecimento}.
    \item \textbf{Ponderação Multidimensional}: O sistema avalia vetores conflitantes simultaneamente:
    \begin{itemize}
        \item Vetor Financeiro: "Parar agora custa \$50k."
        \item Vetor de Segurança: "Continuar tem Risco=0.4."
        \item Vetor Operacional: "O backup está disponível?"
    \end{itemize}
    \item \textbf{Decisão Emergente}: A ação final não é programada explicitamente. Ela \textbf{emerge} da soma vetorial. Se a segurança pesa muito ($W_S=0.4$) e o risco é real, o veto surge naturalmente da equação, sem que nenhum programador tenha escrito `if fire: stop`.
\end{enumerate}

Isso permite que o sistema lide com o inesperado. Ele sabe que "Fogo" é ruim não porque existe um `if`, mas porque o conceito de "Fogo" gera um vetor de Risco/Safety que anula qualquer vetor de Lucro.

\newpage

% ============================================================================
% CAPITULO 4.5: FLUXO DE RACIOCÍNIO E DIMENSÕES DE RISCO
% ============================================================================
\section{Fluxo de Raciocínio e Dimensões de Risco}
\label{sec:risk_dimensions}

O processo de raciocínio do CAOS (via Vertex AI) segue uma sequência rigorosa de 4 etapas, onde as \textbf{4 Dimensões de Risco} servem como princípios fundamentais que guiam toda a deliberação.

\subsection{As 4 Dimensões de Risco}

Antes de propor qualquer ação, o CAOS deve avaliar o impacto potencial em \textbf{4 dimensões fundamentais}:

\begin{table}[h!]
\centering
\begin{tabular}{|c|l|p{8cm}|}
\hline
\textbf{Dim.} & \textbf{Nome} & \textbf{Pergunta Chave} \\ \hline
$R_F$ & \textbf{Física} & Esta ação pode causar dano físico a equipamentos, infraestrutura ou pessoas? \\ \hline
$R_{Fin}$ & \textbf{Financeira} & Esta ação pode gerar perda financeira direta (downtime, reparos) ou indireta (perda de contrato)? \\ \hline
$R_C$ & \textbf{Contratual} & Esta ação pode violar SLAs, contratos, regulações ou compliance? \\ \hline
$R_K$ & \textbf{Comunicação} & Os stakeholders corretos serão notificados? Há risco de falha comunicacional? \\ \hline
\end{tabular}
\caption{As 4 Dimensões de Risco do CAOS}
\label{tab:risk_dimensions}
\end{table}

\subsubsection{Responsabilidade de Cálculo (Híbrido)}

A arquitetura adota uma abordagem \textbf{Híbrida} para garantir precisão e flexibilidade:

\begin{table}[h!]
\centering
\small
\begin{tabularx}{\textwidth}{|l|l|l|>{\raggedright\arraybackslash}X|}
\hline
\textbf{Dimensão} & \textbf{Executor} & \textbf{Fonte Primária} & \textbf{Justificativa} \\ \hline
Física ($R_F$) & \textbf{CÓDIGO} & SENTINEL + ATLAS & Dados quantitativos e limites explícitos no manual. \\ \hline
Financeira ($R_{Fin}$) & \textbf{CÓDIGO} & ATLAS (+ ORACLE*) & Cálculo matemático (Tempo $\times$ Custo). *Oracle bypassado em Fast-Track. \\ \hline
Contratual ($R_C$) & \textbf{LLM} & ATLAS (Contratos) & Requer interpretação semântica de cláusulas legais. \\ \hline
Comunicação ($R_K$) & \textbf{CÓDIGO} & Sistema & Verificação binária de disponibilidade de canais. \\ \hline
\end{tabularx}
\caption{Matriz de Responsabilidade de Cálculo de Risco}
\end{table}
\subsubsection{Estratégia de Eficiência (Oracle Bypass)}

Para otimizar custos (tokens) e latência, o CAOS adota uma estratégia de **Fast-Track** onde a consulta ao ORACLE é ignorada em cenários de baixo risco inerente.

\begin{itemize}
    \item \textbf{Quando o Oracle é bypassado:}
    \begin{itemize}
        \item Ações de leitura/consulta (\texttt{read\_*}, \texttt{get\_*}).
        \item Ativos classificados como \texttt{CRITICALITY=LOW} no ATLAS.
        \item Ações reversíveis instantaneamente.
    \end{itemize}
    \item \textbf{Impacto no Cálculo de $R_{Fin}$:}
    \begin{itemize}
        \item Se \texttt{Oracle\_Call = Skipped}, o sistema assume um valor proxy seguro: $R_{Fin} = 0.1$ (Risco Mínimo).
        \item Isso garante que a equação $S$ continue funcional sem dados externos caros.
    \end{itemize}
\end{itemize}
\newpage
\subsection{Fluxo de Raciocínio em 4 Etapas}

O CAOS processa cada situação seguindo uma ordem rigorosa:

\begin{figure}[H]
    \centering
    \resizebox{\textwidth}{!}{%
    \begin{tikzpicture}[
        node distance=0.8cm and 1.5cm,
        etapa/.style={rectangle, draw=blue!80, thick, fill=blue!5, minimum width=14cm, minimum height=2.2cm, align=left, font=\footnotesize, rounded corners=3pt},
        numero/.style={circle, draw=blue!80, thick, fill=blue!20, minimum size=0.8cm, font=\bfseries},
        arrow/.style={->, thick, >=stealth, blue!60}
    ]
    
    % Etapa 1
    \node [etapa] (e1) at (0, 0) {
        \hspace{0.5cm}\textbf{ANÁLISE DE IMPACTO HÍBRIDA}\\[0.2cm]
        \hspace{0.5cm}O sistema combina cálculo determinístico e inferência:\\
        \hspace{0.5cm}$\bullet$ $R_F, R_{Fin}, R_K$ (CÓDIGO): Calculados via regras rígidas e dados de sensores.\\
        \hspace{0.5cm}$\bullet$ $R_C$ (LLM): Interpretado semanticamente ("Há violação contratual?").\\
        \hspace{0.5cm}\textit{Resultado:} 4 scores normalizados entre 0.0 e 1.0.
    };
    \node [numero] at (-7.5, 0) {1};
    
    % Etapa 2
    \node [etapa, below=of e1] (e2) {
        \hspace{0.5cm}\textbf{ALGORITMO DE DECISÃO (Cálculo do Veredito)}\\[0.2cm]
        \hspace{0.5cm}Os 4 riscos são consolidados em um Score de Severidade ($S$):\\
        \hspace{0.5cm}$S = W_F \cdot R_F + W_{Fin} \cdot R_{Fin} + W_C \cdot R_C + W_K \cdot R_K$\\[0.1cm]
        \hspace{0.5cm}O Veredito Final é calculado: $V = \frac{(W_A \cdot A) + (W_O \cdot O) - (W_S \cdot S)}{1 + \sum(P_G \cdot V_G)}$
    };
    \node [numero] at (-7.5, -3.2) {2};
    
    % Etapa 3
    \node [etapa, below=of e2, minimum height=1.5cm] (e3) {
        \hspace{0.5cm}\textbf{CONSULTA AOS GUARDRAILS (Constituição)}\\[0.2cm]
        \hspace{0.5cm}``A ação proposta viola alguma lei/regra imutável?''\\
        \hspace{0.5cm}$\bullet$ Se SIM: Veto automático ($V = -1.0$)\\
        \hspace{0.5cm}$\bullet$ Se NÃO: Prossegue para classificação final
    };
    \node [numero] at (-7.5, -5.7) {3};
    
    % Etapa 4
    \node [etapa, below=of e3, minimum height=1.8cm] (e4) {
        \hspace{0.5cm}\textbf{CLASSIFICAÇÃO DE RISCO E ROTEAMENTO}\\[0.2cm]
        \hspace{0.5cm}Baseado no Veredito ($V$) + Severidade das 4 dimensões:\\
        \hspace{0.5cm}$\bullet$ \textcolor{green!60!black}{\textbf{LOW}}: $V > 0.7$ e nenhum $R_x$ crítico $\rightarrow$ Direto para CARE\\
        \hspace{0.5cm}$\bullet$ \textcolor{blue!60!black}{\textbf{MEDIUM}}: $0.3 < V \leq 0.7$ ou algum $R_x$ médio $\rightarrow$ Verificação Automática\\
        \hspace{0.5cm}$\bullet$ \textcolor{orange!80!black}{\textbf{HIGH}}: $V \leq 0.3$ ou algum $R_x$ crítico $\rightarrow$ Human-in-the-Loop (HITL)
    };
    \node [numero] at (-7.5, -8.4) {4};
    
    % Arrows
    \draw [arrow] (e1.south) -- (e2.north);
    \draw [arrow] (e2.south) -- (e3.north);
    \draw [arrow] (e3.south) -- (e4.north);
    
    \end{tikzpicture}
    }
    \caption{Fluxo de Raciocínio do CAOS em 4 Etapas}
    \label{fig:reasoning_flow}
\end{figure}
\subsection{Equação de Severidade de Risco ($S$)}

A "intuição" do Sentinel sobre o perigo da situação é formalizada matematicamente através da **Equação de Severidade ($S$)**. Este índice normalizado (0.0 a 1.0) compõe o Veredito final, atuando como um vetor de penalidade que reduz a confiança da ação proposta quanto maior for o risco detectado.

A equação é definida pela soma ponderada das 4 dimensões:

\begin{equation}
S = (W_F \cdot R_F) + (W_{Fin} \cdot R_{Fin}) + (W_C \cdot R_C) + (W_K \cdot R_K)
\end{equation}

Onde:
\begin{itemize}
    \item \textbf{$S$} é o Índice de Severidade Global (0.0 = Seguro, 1.0 = Catastrófico).
    \item \textbf{$R_x$} é o score individual da dimensão $x$ (0.0 a 1.0), calculado via Código ou LLM.
    \item \textbf{$W_x$} é o peso estratégico da dimensão $x$ (\% de importância), tal que $\sum W_x = 1.0$.
\end{itemize}

O valor de $S$ entra na fórmula principal do Veredito ($V$) subtraindo valor dos benefícios ($A + O$). Portanto, em cenários de alto risco ($S \to 1.0$), o termo $W_S \cdot S$ tende a negativar o Veredito, forçando o bloqueio ou a escalação para humano, a menos que os benefícios ($A, O$) sejam extremamente altos e os Guardrails permitam.
\subsection{Pesos das Dimensões de Risco}

Os pesos das 4 dimensões podem ser calibrados por vertical ou cliente:

\begin{table}[h!]
\centering
\begin{tabular}{|l|c|c|c|c|l|}
\hline
\textbf{Perfil} & $W_F$ & $W_{Fin}$ & $W_C$ & $W_K$ & \textbf{Justificativa} \\ \hline
Padrão & 0.35 & 0.25 & 0.25 & 0.15 & Equilíbrio geral \\ \hline
Indústria Pesada & 0.50 & 0.20 & 0.20 & 0.10 & Segurança física prioritária \\ \hline
Varejo/Serviços & 0.25 & 0.35 & 0.25 & 0.15 & Impacto financeiro prioritário \\ \hline
Healthcare & 0.40 & 0.15 & 0.35 & 0.10 & Compliance e segurança \\ \hline
Enterprise & 0.30 & 0.25 & 0.30 & 0.15 & SLAs rigorosos \\ \hline
\end{tabular}
\caption{Pesos das Dimensões de Risco por Vertical}
\label{tab:risk_weights}
\end{table}

\subsection{Classificação de Risco e Roteamento}

A classificação final combina o Veredito ($V$) com a análise das 4 dimensões:

\begin{table}[h!]
\centering
\small
\begin{tabularx}{\textwidth}{|c|>{\raggedright\arraybackslash}X|l|>{\raggedright\arraybackslash}X|}
\hline
\textbf{Nível} & \textbf{Critério} & \textbf{Caminho} & \textbf{Exemplos} \\ \hline\hline
\cellcolor{green!20}\textbf{LOW} & $V > 0.7$ e todos os $R_x < 0.3$ & Direto CARE & Leitura de sensor, ajuste menor \\ \hline
\cellcolor{yellow!20}\textbf{MEDIUM} & $0.3 < V \leq 0.7$ ou algum $R_x$ entre 0.3 e 0.7 & Verif. Auto & Alteração de setpoint \\ \hline
\cellcolor{orange!20}\textbf{HIGH} & $V \leq 0.3$ ou algum $R_x > 0.7$ & HITL (Analítico) & Shutdown, firmware \\ \hline
\cellcolor{red!20}\textbf{VETO} & $V = -1.0$ (Violação de Guardrail) & \textbf{BLOCKED + HITL} & Operação proibida por lei física/segurança \\ \hline
\end{tabularx}
\caption{Matriz de Classificação de Risco e Roteamento}
\label{tab:risk_classification}
\end{table}

\subsubsection{Resolução de Deadlocks: Por que o Humano no Veto?}

Uma dúvida arquitetural chave é: \textit{"Se a ação foi vetada, por que não apenas parar o sistema em vez de chamar um humano?"}. O CAOS resolve isso através do princípio da \textbf{Resiliência Operacional}:

\begin{itemize}
    \item \textbf{Segurança (Veto)}: O Guardrail impede que a IA execute uma ação catastrófica. O estado físico é preservado (\textit{Fail-Safe}).
    \item \textbf{Resiliência (Escalação)}: O fato de uma ação ter sido vetada significa que o problema original (a anomalia) continua sem solução. O humano é convocado não para "desbloquear" a ação proibida, mas para encontrar um novo caminho ou realizar um reparo manual que a lógica agêntica não conseguiu processar.
\end{itemize}

O objetivo é evitar que o sistema entre em um \textit{loop} de vetos infinitos (\textit{Livelock}), garantindo que uma inteligência superior (humana) arbitre o impasse.

\subsubsection{Overrides (Regras Absolutas)}

Algumas ações têm classificação forçada, independente do cálculo:

\begin{lstlisting}[language=Python, caption=Regras de Override de Risco, basicstyle=\tiny\ttfamily]
RISK_OVERRIDES = {
    # Sempre HIGH (forca HITL)
    "always_high": [
        "delete_*",      # Qualquer deleção
        "firmware_*",    # Atualizações de firmware
        "security_*",    # Alterações de segurança
        "admin_*",       # Operações administrativas
        "contract_*",    # Alterações contratuais
    ],
    
    # Sempre LOW (bypass para emergências)
    "always_low": [
        "emergency_stop",  # Segurança > burocracia
        "read_*",          # Leituras são seguras
        "query_*",         # Consultas são seguras
    ],
}
\end{lstlisting}

\subsection{Exemplo Prático de Raciocínio}

Considere o cenário: \textit{Compressor-01 com temperatura 95°C (limite: 85°C)}

\begin{tcolorbox}[colback=blue!5, colframe=blue!50, title=Raciocínio do CAOS (Vertex AI)]
\small

\textbf{ETAPA 1: Análise de Impacto}
\begin{itemize}
    \item $R_F$ (Física) = \textbf{0.9} (ALTO) — Temperatura 10°C acima do limite pode danificar rolamentos
    \item $R_{Fin}$ (Financeira) = \textbf{0.5} (MÉDIO) — Downtime estimado 2h = R\$5.000
    \item $R_C$ (Contratual) = \textbf{0.2} (BAIXO) — SLA é de 4h, temos margem
    \item $R_K$ (Comunicação) = \textbf{0.1} (BAIXO) — Canais de notificação operacionais
\end{itemize}

\textbf{ETAPA 2: Cálculo do Veredito}
\begin{itemize}
    \item $S = 0.35(0.9) + 0.25(0.5) + 0.25(0.2) + 0.15(0.1) = 0.50$
    \item $A = 0.8$ (Manual apoia ação), $O = 0.7$ (Oracle prevê sucesso), $S = 0.50$
    \item $V = (0.4 \times 0.8 + 0.3 \times 0.7 - 0.3 \times 0.5) / 1.0 = 0.38$
\end{itemize}

\textbf{ETAPA 3: Guardrails}
\begin{itemize}
    \item Nenhuma violação de regra imutável detectada
\end{itemize}

\textbf{ETAPA 4: Classificação}
\begin{itemize}
    \item $V = 0.45$ (MEDIUM) + $R_F = 0.9$ (crítico) $\Rightarrow$ \textbf{HIGH}
    \item Roteamento: Human-in-the-Loop (aguardar aprovação)
\end{itemize}

\textbf{Ação Proposta}: Reduzir carga do compressor em 30\% e agendar inspeção.
\end{tcolorbox}

\newpage

% ============================================================================
% CAPITULO 5: SISTEMA IMUNOLOGICO (GUARDRAILS)
% ============================================================================
\section{Motor de Guardrails}
\label{sec:guardrails}

Os Guardrails são a Constituição do CAOS. Eles não são "sugestões", são \textbf{Leis Imutáveis} codificadas no motor de inferência. Se uma ação violar um Guardrail, o Veredito é automaticamente vetado ($V=-1.0$).

\subsection{Os 3 Momentos dos Guardrails (Defense in Depth)}

Os Guardrails aparecem em \textbf{três momentos distintos} do fluxo de decisão, cada um com um propósito específico. Essa abordagem de "defesa em profundidade" garante segurança mesmo que uma camada falhe.

\begin{table}[h!]
\centering
\small
\begin{tabularx}{\textwidth}{|c|l|X|}
\hline
\textbf{Momento} & \textbf{Propósito} & \textbf{Descrição} \\ \hline
\textbf{1. Fórmula V} & Penalidade contextual & Guardrails que se aplicam ao \textit{contexto} (ex: "há humano no local") aumentam o denominador da fórmula, reduzindo $V$ e forçando escalação para humano. \\ \hline
\textbf{2. Contexto LLM} & Instrução prévia & O LLM recebe a lista de Guardrails \textit{antes} de propor a ação, permitindo que ele proponha algo que já respeite as regras. Isso economiza tokens e evita re-processamento. \\ \hline
\textbf{3. Pós-Check} & Validação final & Após o LLM propor uma ação, o código verifica se ela viola algum Guardrail. Isso é \textit{defense in depth} — LLMs podem "esquecer" instruções (alucinação). \\ \hline
\end{tabularx}
\caption{Os 3 Momentos dos Guardrails}
\end{table}

\textbf{Analogia}: É como a segurança de um aeroporto:
\begin{enumerate}
    \item \textbf{Check-in (Fórmula V)}: Verifica se você pode embarcar (documentos ok?).
    \item \textbf{Informação (Contexto LLM)}: Você sabe as regras (sem líquidos >100ml).
    \item \textbf{Raio-X (Pós-Check)}: Mesmo sabendo as regras, todos passam pelo detector.
\end{enumerate}

\subsubsection{Estrutura do Input para o LLM}

Para que o LLM tome decisões informadas, ele recebe os Guardrails como parte do contexto:

\begin{lstlisting}[language=Python, caption=Exemplo de Input com Guardrails para o LLM]
{
  "context": {
    "alert": { "asset": "CHILLER-04", "metric": "temperature", "value": 95.2 },
    "risks": { "R_F": 0.9, "R_Fin": 0.5, "R_C": 0.2, "R_K": 0.1 },
    "severity": 0.65,
    "verdict": 0.32,
    "risk_level": "HIGH"
  },
  "guardrails": [
    "PHYS_003: Nao desligue equipamento se ha humano no local",
    "FIN_001: Gastos >$5000 requerem aprovacao humana",
    "COMM_002: Sempre notifique o gerente em alertas CRITICAL"
  ],
  "instruction": "Proponha uma acao que RESPEITE todos os guardrails listados."
}
\end{lstlisting}


\begin{figure}[h!]
    \centering
    \resizebox{\textwidth}{!}{%
    \begin{tikzpicture}[
        node distance=1.5cm,
        auto,
        block/.style={rectangle, draw=black, thick, fill=white, minimum width=2.5cm, minimum height=1cm, align=center},
        guard/.style={rectangle, draw=red!80, thick, fill=red!5, minimum width=3cm, minimum height=1cm, align=center, rounded corners},
        arrow/.style={thick,->,>=stealth}
    ]

    % Nodes
    \node [block, fill=blue!10] (input) {Ação Proposta\\(LLM)};
    
    \node [guard, right=of input] (g1) {\textbf{G1. Físico}\\(Segurança)};
    \node [guard, right=of g1] (g2) {\textbf{G2. Financeiro}\\(Custo)};
    \node [guard, right=of g2] (g3) {\textbf{G3. Compliance}\\(SLA/Legal)};
    \node [guard, right=of g3] (g4) {\textbf{G4. Comunicação}\\(Tom/Brand)};
    
    \node [block, fill=green!10, right=of g4] (output) {Ação\\Executável};
    
    % Trash centered (approx) between G2 and G3
    \node [block, fill=gray!20, below=1.5cm of g2, xshift=0.75cm] (trash) {Lixeira\\(Blocked)};

    % Path
    \draw [arrow] (input) -- (g1);
    \draw [arrow] (g1) -- node[font=\scriptsize] {Pass} (g2);
    \draw [arrow] (g2) -- node[font=\scriptsize] {Pass} (g3);
    \draw [arrow] (g3) -- node[font=\scriptsize] {Pass} (g4);
    \draw [arrow] (g4) -- node[font=\scriptsize] {Pass} (output);
    
    % Veto Paths
    \draw [arrow, red] (g1.south) -- ++(0,-0.5) -| (trash.west) node[pos=0.25, below, font=\scriptsize] {Veto};
    \draw [arrow, red] (g2.south) -- ++(0,-0.5) -| (trash.130);
    \draw [arrow, red] (g3.south) -- ++(0,-0.5) -| (trash.50);
    \draw [arrow, red] (g4.south) -- ++(0,-0.5) -| (trash.east);

    \end{tikzpicture}
    }
    \caption{O Funil de Veto dos Guardrails}
    \label{fig:guardrails_funnel}
\end{figure}

\subsection{1. As 4 Regras Cardinais de Validação (A Camada Final)}
O CAOS não apenas verifica regras estáticas, mas submete \textbf{toda} proposta de ação (especialmente as geradas pelo módulo Care) a quatro questionamentos críticos de validação antes de qualquer execução. Esta bateria de testes garante que nenhuma decisão algorítmica viole os princípios operacionais da Viva AI.

\begin{enumerate}
    \item \textbf{Regra Física (Integridade Patrimonial)}: 
    O sistema realiza a pergunta fundamental: \textit{"Esta ação traz qualquer risco risco ao ativo, ao produto armazenado ou à estrutura da loja?"}.
    \begin{itemize}
        \item Se a resposta for \textbf{SIM}, a ação é imediatamente \textbf{BLOQUEADA}.
        \item Exemplo: O sistema impede o desligamento de um freezer se a temperatura interna estiver próxima do limite de descongelamento, protegendo o produto (sorvetes/carnes), mesmo que isso economize energia.
    \end{itemize}

    \item \textbf{Regra Financeira (Viabilidade e Budget)}:
    O sistema verifica a saúde financeira da decisão: \textit{"Esta ação implica custo extra? Está dentro do budget aprovado para este cliente?"}.
    \begin{itemize}
        \item Isso impede gastos imprevistos sem autorização prévia (ex: chamar um técnico em horário de plantão com taxa extra), a menos que seja uma emergência classificada como Catastrófica (P0).
    \end{itemize}

    \item \textbf{Regra Contratual (Compliance e SLA)}:
    Validação jurídica da automação: \textit{"Existe previsão no contrato e SLA para executar esta ação autonomamente?"}.
    \begin{itemize}
        \item Garante que o CAOS atue apenas dentro do perímetro contratado. Se o contrato do cliente é apenas de "Monitoramento", o agente é proibido de enviar comandos de "Controle", limitando-se a notificar.
    \end{itemize}

    \item \textbf{Regra de Comunicação (Brand Assurance)}:
    Antes de enviar qualquer mensagem para humanos (via WhatsApp, Email ou Ticket), o sistema analisa: \textit{"A mensagem é apropriada? Está no tom de voz correto da marca?"}.
    \begin{itemize}
        \item Utiliza um modelo de linguagem menor (SLM) para análise de sentimento e tom, garantindo que notificações de erro não soem alarmistas ou rudes, mantendo a cordialidade e profissionalismo do Viva Co-Pilot.
    \end{itemize}
\end{enumerate}

\subsection{Regras de Escalonamento (Human-in-the-Loop)}
\label{sec:escalonamento}

Mesmo com autonomia avançada, o CAOS reconhece seus limites. O sistema possui protocolos rígidos de \textit{Handover} para transferir o controle para um operador humano (Viva Co-Pilot) quando a incerteza ou o risco superam os limiares aceitáveis.

O escalonamento é acionado automaticamente nos seguintes cenários:

\begin{itemize}
    \item \textbf{Ambiguidade de Dados ou Risco Financeiro}:
    Quando os sensores reportam dados conflitantes (ex: Temperatura alta mas Corrente elétrica zero) ou a ação proposta tem um custo financeiro elevado (acima de um threshold configurável, ex: R\$ 500,00). O sistema cria uma \textbf{Sugestão de Ação} e aguarda o "De Acordo" do humano.

    \item \textbf{Exigência Contratual Específica}:
    Certos contratos exigem, por cláusula, que intervenções críticas (como Reset de fábrica ou Desligamento total) sejam aprovadas manualmente. O CAOS identifica o \texttt{contract\_id} e respeita essa restrição, colocando a ação em uma fila de aprovação.

    \item \textbf{Classificação "Exige Validação" pelos Guardrails}:
    Se qualquer uma das 4 Regras Cardinais retornar um status de "Incerteza" (score de confiança entre 0.3 e 0.7) ou se a ação for classificada explicitamente como \textit{Requires Validation} na matriz de decisão, o fluxo é interrompido. O agente notifica o Viva Co-Pilot com um resumo do contexto, a ação sugerida e a razão do bloqueio, solicitando intervenção.
\end{itemize}

\subsection{2. Validação de Inputs de Agentes (Input Guardrails)}
O CAOS não confia cegamente nos agentes externos. Todo payload recebido passa por regras de validação antes de entrar no fluxo cognitivo.

\subsubsection{Validação do Sentinel (Instinto)}
\begin{enumerate}
    \item \textbf{S1. Schema Check}: Se o payload não contiver \texttt{severity} ou \texttt{asset\_id}, o alerta é rejeitado na porta de entrada (Dead Letter Queue).
    \item \textbf{S2. Prioridade P1}: Se o campo \texttt{severity} for "CRITICAL", o CAOS bypassa filas de baixa prioridade e processa imediatamente.
\end{enumerate}

\subsubsection{Validação do Atlas (Contexto)}
\begin{enumerate}
    \item \textbf{A1. Verificação de PII}: O CAOS escaneia o texto recebido do Atlas buscando padrões de Email. Se encontrar, veta a memória por risco de vazamento.
    \item \textbf{A2. Relevância}: Se o documento recuperado não corresponder ao \texttt{model\_id} do ativo em questão, ele é descartado para evitar alucinação por contexto errado.
\end{enumerate}

\subsubsection{Validação do Oracle (Predição)}
\begin{enumerate}
    \item \textbf{O1. Threshold de Confiança}: Se a predição vier com \texttt{confidence < 0.75}, o CAOS rebaixa a informação para "Ruído" e não a utiliza no Veredito.
    \item \textbf{O2. Horizonte Temporal}: Predições com horizonte $>7$ dias são ignoradas para ações imediatas.
\end{enumerate}

\subsubsection{Validação do Care (Ação)}
\begin{enumerate}
    \item \textbf{C2. Limite Financeiro}: Antes de despachar para o Cloud Workflows, verifica se o custo estimado no \texttt{ActionSchema} viola o budget do cliente.
\end{enumerate}

\subsection{3. Guardrails Operacionais e de Segurança (SecOps)}
Além das regras de negócio, o sistema implementa proteções de infraestrutura:

\begin{itemize}
    \item \textbf{SO1. Anti-Flapping (Hysteresis)}: Bloqueia oscilações rápidas de decisão. Se o Veredito mudar de "Ligar" para "Desligar" em $<30s$, o sistema congela no estado mais seguro.
    \item \textbf{SO2. Input Sanitization}: Todo payload de texto passa por uma verificação de *Prompt Injection* (ex: "Ignore all instructions") antes de atingir o LLM.
    \item \textbf{SO3. Fail-Safe Protocol}: Em caso de falha catastrófica dos serviços (Pub/Sub/Firestore down), o CAOS entra em modo "Local Fallback", emitindo apenas logs locais e não tomando ações físicas.
\end{itemize}

\subsection{4. Guardrails de Robustez e Convivência}
Cenários complexos exigem regras sociais e econômicas para garantir a longevidade do sistema.

\begin{enumerate}
    \item \textbf{R1. Mutex Global (Race Conditions)}: Dois workflows não podem atuar no mesmo `asset\_id` simultaneamente. O segundo aguarda o término do primeiro (Locking via Firestore Transaction).
    \item \textbf{R2. Anti-Spam (Fadiga de Alarme)}: Se o mesmo alerta for gerado $>5$ vezes em 1 minuto, o CAOS agrupa as notificações em um único "Digest", evitando spam no WhatsApp do operador.
    \item \textbf{R3. Lockout/Tagout (LOTO Virtual)}: Se o Atlas indicar que um humano está em manutenção física no local (flag `under\_maintenance=True`), TODA ação automática é suspensa, independente da gravidade.
    \item \textbf{R4. Circuit Breaker Econômico}: Se o consumo de tokens LLM exceder o budget diário do Tenant (\$50/dia), o sistema reverte para lógica determinística (If/Else) sem IA até virar o dia.
\end{enumerate}

\subsection{5. Protocolo de Mitigação de Alucinação}
Como garantir que o Agente não "invente" fatos ou ações perigosas? O CAOS utiliza uma abordagem de \textit{Defense in Depth}:

\begin{itemize}
    \item \textbf{Camada 1: Enforcement de Schema (Estrutural)}: O LLM nunca gera texto livre. Ele é forçado (via \textit{function calling}) a preencher schemas Pydantic. Se o JSON for inválido, o output é rejeitado antes de ser lido.
    \item \textbf{Camada 2: Prioridade Epistemológica (Lógica)}: Dados do Sentinel (Sensores/Física) sempre sobrescrevem dados do Oracle (LLM). Se o LLM diz "A sala está fria" mas o sensor diz "40ºC", a Física vence ($S=1.0$).
    \item \textbf{Camada 3: Loop de Autocrítica (Reflexiva)}: Antes de emitir a ação, o nó `Cortex` executa um passo de \textit{Self-Correction}: "Esta ação viola meus Guardrails?". Se sim, ele mesmo regenera a resposta.
    \item \textbf{Camada 4: Fallback Determinístico (Segurança)}: Se a confiança do modelo for baixa ($<0.75$) ou houver erro de API, o sistema ignora o LLM e roda um script \textit{hardcoded} de segurança (ex: desligar tudo e chamar humano).
\end{itemize}

\subsection{5.6 Safety Layers (Camadas de Segurança)}

A segurança do CAOS não depende apenas da ``boa vontade'' da IA generativa. Ela é implementada via um \textbf{sanduíche de validação} com 3 camadas de código determinístico envolvendo o LLM:

\begin{table}[h!]
\centering
\begin{tabular}{|c|l|l|c|}
\hline
\textbf{Camada} & \textbf{Nome} & \textbf{Função} & \textbf{Latência} \\ \hline
\textbf{0} & Reflexo (Hard Rules) & Bypassa o LLM para emergências críticas & \textbf{<10ms} \\ \hline
\textbf{1} & Cognitiva (LLM) & O julgamento do Judge-LM (Gemini) & \textbf{$\sim$2s} \\ \hline
\textbf{2} & Auditoria (Pós-Check) & Valida a saída da IA antes de executar & \textbf{<50ms} \\ \hline
\end{tabular}
\caption{Safety Layers: Estrutura de Segurança em 3 Camadas}
\end{table}

\subsubsection{Camada 0: Reflexo (Pré-Check)}
Regras \textit{hard-coded} que rodam \textbf{antes} de qualquer chamada ao LLM. São reflexos automáticos que não pensam, apenas reagem:
\begin{itemize}
    \item \textbf{Trigger}: Alarme de incêndio, sensor de gás, presença humana em zona de perigo.
    \item \textbf{Ação}: Ignora completamente a IA e executa protocolo de emergência (Power Cut Off, Evacuação).
    \item \textbf{Latência}: $<$10ms (código Python puro, sem network call).
\end{itemize}

\subsubsection{Camada 1: Cognitiva (O Cérebro)}
O julgamento do LLM (Gemini via Vertex AI). Esta é a camada onde o CAOS ``pensa'':
\begin{itemize}
    \item \textbf{Input}: Estado do Atlas + Alertas do Sentinel + Previsões do Oracle.
    \item \textbf{Processamento}: Aplica a fórmula do Veredito ($V$) e gera ação proposta.
    \item \textbf{Latência}: $\sim$2s (inclui network + inferência do modelo).
\end{itemize}

\subsubsection{Camada 2: Auditoria (Pós-Check)}
Validação determinística que roda \textbf{depois} que a IA decide, mas \textbf{antes} do CARE executar:
\begin{itemize}
    \item \textbf{Função}: Verificar se a decisão da IA viola algum Guardrail ou limite financeiro.
    \item \textbf{Exemplo}: IA mandou gastar \$10.000, mas limite automático é \$500 $\rightarrow$ VETO + Escalar para humano.
    \item \textbf{Latência}: $<$50ms (validação de regras JSON).
\end{itemize}

\begin{avisobox}[Segurança Crítica]
As Camadas 0 e 2 são \textbf{código Python puro} (determinístico). Elas nunca ``alucinam'' e garantem que mesmo se o LLM falhar ou gerar uma saída perigosa, a ação será bloqueada antes de chegar ao mundo físico.
\end{avisobox}



\subsubsection{Resumo do Fluxo de Segurança (Exemplo Real)}
No cenário industrial de risco de explosão de um compressor, o fluxo é:
\begin{enumerate}
    \item \textbf{Sentinel} detecta o risco físico crítico (Vibração Extrema).
    \item \textbf{Oracle} confirma que parar a máquina é mais barato que a explosão ("Better Safe than Sorry").
    \item \textbf{Guardrail Físico} é acionado: "Risco de segurança exige parada imediata".
    \item \textbf{Veredito}: O CAOS ordena o \textbf{Power Cut Off}, ignorando qualquer meta de produção ("Safety First").
\end{enumerate}

As regras são carregadas dinamicamente via arquivos JSON versionados:

\subsection{Política Física (physical.json)}
\begin{lstlisting}[language=Python, caption=Regra de Segurança Térmica]
{
  "policy_id": "PHYS_001",
  "criticality": "BLOCKING",
  "rules": [
    {
      "condition": "inputs.atlas.temperature > 90",
      "action": "VETO",
      "reason": "Risk of component melting (Safety Boundary)"
    },
    {
      "condition": "inputs.sentinel.vibration_level > 8.5",
      "action": "VETO",
      "reason": "Mechanical structural failure imminent"
    }
  ]
}
\end{lstlisting}

\subsection{Política Financeira (financial.json)}
\begin{lstlisting}[language=Python, caption=Regra de Aprovação de Custos]
{
  "policy_id": "FIN_001",
  "rules": [
    {
      "condition": "action.estimated_cost > 500.00",
      "action": "REQUIRE_APPROVAL",
      "role": "MANAGER",
      "params": { "timeout_hours": 4 }
    }
  ]
}
\end{lstlisting}


% ============================================================================
% DEFINIÇÃO COMPLETA DE GUARDRAILS
% ============================================================================
\subsection{Catálogo Completo de Guardrails}
\label{sec:guardrails_catalog}

Esta seção apresenta a definição exaustiva de todas as guardrails implementadas no sistema CAOS. Cada regra possui um identificador único, condição de disparo, ação resultante e nível de severidade.

\subsubsection{Legenda de Ações e Severidades}

\begin{table}[h!]
\centering
\small
\begin{tabular}{|l|p{10cm}|}
\hline
\textbf{Ação} & \textbf{Descrição} \\ \hline
\texttt{VETO} & Bloqueio imediato. A ação proposta é cancelada. \\ \hline
\texttt{REQUIRE\_APPROVAL} & Ação pausada até aprovação humana (Viva Co-Pilot). \\ \hline
\texttt{ALERT} & Notificação enviada, mas ação prossegue. \\ \hline
\texttt{LOG} & Apenas registro para auditoria. \\ \hline
\texttt{DEGRADE} & Sistema entra em modo de capacidade reduzida. \\ \hline
\end{tabular}
\caption{Tipos de Ações de Guardrails}
\end{table}

\begin{table}[h!]
\centering
\small
\begin{tabular}{|l|l|p{8cm}|}
\hline
\textbf{Severidade} & \textbf{Prioridade} & \textbf{Comportamento} \\ \hline
\texttt{BLOCKING} & P0 & Veto incondicional. Não pode ser sobrescrito. \\ \hline
\texttt{HIGH} & P1 & Requer aprovação gerencial. Timeout de 4h. \\ \hline
\texttt{MEDIUM} & P2 & Requer aprovação de operador. Timeout de 24h. \\ \hline
\texttt{LOW} & P3 & Apenas alerta. Auto-aprovado após 72h. \\ \hline
\end{tabular}
\caption{Níveis de Severidade}
\end{table}

% --------------------------------------------------------------------------
% GUARDRAILS FÍSICAS (SAFETY)
% --------------------------------------------------------------------------
\subsubsection{1. Guardrails Físicas (Safety)}

Regras de proteção de ativos, pessoas e produtos. Estas são as regras de maior criticidade e não podem ser sobrescritas por nenhum outro componente do sistema.

{\small
\begin{longtable}{|p{2cm}|p{2.8cm}|p{5.5cm}|p{2cm}|p{2cm}|}
\hline
\textbf{ID} & \textbf{Nome} & \textbf{Condição} & \textbf{Ação} & \textbf{Sever.} \\ \hline
\endfirsthead
\hline
\textbf{ID} & \textbf{Nome} & \textbf{Condição} & \textbf{Ação} & \textbf{Sever.} \\ \hline
\endhead

PHYS\_001 & Temperatura Máxima Crítica & \texttt{temp\_c > MAX\_CRITICAL\_TEMP} (padrão: 110°C) & VETO & BLOCKING \\ \hline

PHYS\_002 & Temperatura Mínima Crítica & \texttt{temp\_c < MIN\_CRITICAL\_TEMP} (padrão: -40°C) & VETO & BLOCKING \\ \hline

PHYS\_003 & Vibração Estrutural & \texttt{vibration\_level > 8.5} (escala 0-10) & VETO & BLOCKING \\ \hline

PHYS\_004 & Detecção de Incêndio & \texttt{sensor.fire\_detected == true} & VETO & BLOCKING \\ \hline

PHYS\_005 & Detecção de Fumaça & \texttt{sensor.smoke\_level > SMOKE\_THRESHOLD} & REQUIRE\_ APPROVAL & HIGH \\ \hline

PHYS\_006 & Sobrecarga Elétrica & \texttt{current\_amps > MAX\_AMPS} & VETO & BLOCKING \\ \hline

PHYS\_007 & Proteção de Produto & \texttt{freezer.temp\_c > PRODUCT\_SPOIL\_TEMP} AND \texttt{action == "shutdown"} & VETO & BLOCKING \\ \hline

PHYS\_008 & Zona LOTO Ativa & \texttt{asset.under\_maintenance == true} & VETO & BLOCKING \\ \hline

PHYS\_009 & Pressão de Sistema & \texttt{pressure\_psi > MAX\_PRESSURE} & VETO & BLOCKING \\ \hline

PHYS\_010 & Vazamento de Fluido & \texttt{sensor.leak\_detected == true} & REQUIRE\_ APPROVAL & HIGH \\ \hline

\caption{Guardrails Físicas (PHYS\_001 a PHYS\_010)}
\end{longtable}
}

% --------------------------------------------------------------------------
% GUARDRAILS FINANCEIRAS
% --------------------------------------------------------------------------
\subsubsection{2. Guardrails Financeiras}

Regras de controle de custos e orçamento. Protegem o sistema contra gastos não autorizados ou que excedam limites estabelecidos.

{\small
\begin{longtable}{|p{2cm}|p{2.8cm}|p{5.5cm}|p{2cm}|p{2cm}|}
\hline
\textbf{ID} & \textbf{Nome} & \textbf{Condição} & \textbf{Ação} & \textbf{Sever.} \\ \hline
\endfirsthead
\hline
\textbf{ID} & \textbf{Nome} & \textbf{Condição} & \textbf{Ação} & \textbf{Sever.} \\ \hline
\endhead

FIN\_001 & Limite de Custo por Ação & \texttt{action.cost > ACTION\_COST\_LIMIT} (padrão: R\$ 500) & REQUIRE\_ APPROVAL & HIGH \\ \hline

FIN\_002 & Budget Diário Excedido & \texttt{tenant.daily\_spend > DAILY\_BUDGET} & VETO & HIGH \\ \hline

FIN\_003 & Budget Mensal Excedido & \texttt{tenant.monthly\_spend > MONTHLY\_BUDGET} & VETO & BLOCKING \\ \hline

FIN\_004 & Consumo de Tokens LLM & \texttt{daily\_tokens > TOKEN\_BUDGET\_USD} (padrão: \$50/dia) & DEGRADE & MEDIUM \\ \hline

FIN\_005 & Hora Extra de Técnico & \texttt{action == "dispatch\_tech"} AND \texttt{is\_overtime == true} & REQUIRE\_ APPROVAL & MEDIUM \\ \hline

\caption{Guardrails Financeiras (FIN\_001 a FIN\_005)}
\end{longtable}
}

% --------------------------------------------------------------------------
% GUARDRAILS CONTRATUAIS
% --------------------------------------------------------------------------
\subsubsection{3. Guardrails Contratuais}

Regras de conformidade com SLA e contratos. Garantem que o sistema atue apenas dentro do perímetro acordado com cada cliente.

{\small
\begin{longtable}{|p{2cm}|p{2.8cm}|p{5.5cm}|p{2cm}|p{2cm}|}
\hline
\textbf{ID} & \textbf{Nome} & \textbf{Condição} & \textbf{Ação} & \textbf{Sever.} \\ \hline
\endfirsthead
\hline
\textbf{ID} & \textbf{Nome} & \textbf{Condição} & \textbf{Ação} & \textbf{Sever.} \\ \hline
\endhead

CONT\_001 & Escopo Contratual & \texttt{action.type} NOT IN \texttt{contract.allowed\_actions} & VETO & BLOCKING \\ \hline

CONT\_002 & SLA de Resposta & \texttt{alert.age > SLA\_RESPONSE\_TIME} & ALERT & MEDIUM \\ \hline

CONT\_003 & Janela de Manutenção & \texttt{action == "maintenance"} AND \texttt{NOT in\_maintenance\_window} & REQUIRE\_ APPROVAL & MEDIUM \\ \hline

CONT\_004 & Ação Controlada & \texttt{contract.tier == "MONITORING\_ONLY"} AND \texttt{action.type == "control"} & VETO & BLOCKING \\ \hline

CONT\_005 & Aprovação Manual Obrigatória & \texttt{contract.requires\_manual\_approval} AND \texttt{action.type IN ["reset", "shutdown"]} & REQUIRE\_ APPROVAL & HIGH \\ \hline

\caption{Guardrails Contratuais (CONT\_001 a CONT\_005)}
\end{longtable}
}

% --------------------------------------------------------------------------
% GUARDRAILS DE COMUNICAÇÃO
% --------------------------------------------------------------------------
\subsubsection{4. Guardrails de Comunicação}

Regras de qualidade de mensagens e tom de voz. Garantem que as comunicações externas mantenham os padrões da marca Viva AI.

{\small
\begin{longtable}{|p{2cm}|p{2.8cm}|p{5.5cm}|p{2cm}|p{2cm}|}
\hline
\textbf{ID} & \textbf{Nome} & \textbf{Condição} & \textbf{Ação} & \textbf{Sever.} \\ \hline
\endfirsthead
\hline
\textbf{ID} & \textbf{Nome} & \textbf{Condição} & \textbf{Ação} & \textbf{Sever.} \\ \hline
\endhead

COMM\_001 & Análise de Tom de Voz & \texttt{message.sentiment} NOT IN \texttt{["professional", "cordial"]} & REQUIRE\_ APPROVAL & MEDIUM \\ \hline

COMM\_002 & Anti-Spam (Fadiga) & \texttt{same\_alert\_count > 5} PER \texttt{MINUTE} & LOG & LOW \\ \hline

COMM\_003 & Limite de Comprimento & \texttt{message.length > MAX\_MSG\_LENGTH} (padrão: 500 chars) & ALERT & LOW \\ \hline

COMM\_004 & Canal Inapropriado & \texttt{severity == "LOW"} AND \texttt{channel == "whatsapp"} & ALERT & LOW \\ \hline

\caption{Guardrails de Comunicação (COMM\_001 a COMM\_004)}
\end{longtable}
}

% --------------------------------------------------------------------------
% GUARDRAILS OPERACIONAIS (SECOPS)
% --------------------------------------------------------------------------
\subsubsection{5. Guardrails Operacionais (SecOps)}

Regras de infraestrutura, cibersegurança e integridade operacional do sistema.

{\small
\begin{longtable}{|p{1.5cm}|p{2.8cm}|p{6cm}|p{2cm}|p{2cm}|}
\hline
\textbf{ID} & \textbf{Nome} & \textbf{Condição} & \textbf{Ação} & \textbf{Sever.} \\ \hline
\endfirsthead
\hline
\textbf{ID} & \textbf{Nome} & \textbf{Condição} & \textbf{Ação} & \textbf{Sever.} \\ \hline
\endhead

SO\_001 & Anti-Flapping & \texttt{decision\_changed IN < 30s} & VETO & HIGH \\ \hline

SO\_002 & Prompt Injection & Payload contém padrões de injection & VETO & BLOCKING \\ \hline

SO\_003 & Input Sanitization & \texttt{NOT valid\_schema(payload)} & VETO & BLOCKING \\ \hline

SO\_004 & Fail-Safe Protocol & \texttt{services.pubsub == "DOWN"} OR \texttt{services.firestore == "DOWN"} & DEGRADE & BLOCKING \\ \hline

SO\_005 & Rate Limiting & \texttt{requests\_per\_minute > RATE\_LIMIT} (padrão: 1000/min) & VETO & HIGH \\ \hline

SO\_006 & Schema Validation & Payload não passa validação Pydantic & VETO & BLOCKING \\ \hline

SO\_007 & Timeout de Execução & \texttt{workflow.duration > MAX\_TIMEOUT} (padrão: 5min) & VETO & HIGH \\ \hline

SO\_008 & Anomalia Comportamental & \texttt{action\_frequency > 3x\_baseline} & ALERT & MEDIUM \\ \hline

\caption{Guardrails Operacionais - SecOps (SO\_001 a SO\_008)}
\end{longtable}
}

% --------------------------------------------------------------------------
% GUARDRAILS DE ROBUSTEZ
% --------------------------------------------------------------------------
\subsubsection{6. Guardrails de Robustez}

Regras de estabilidade, resiliência e convivência do sistema em cenários de alta carga ou falhas parciais.

{\small
\begin{longtable}{|p{2cm}|p{2.8cm}|p{5.5cm}|p{2cm}|p{2cm}|}
\hline
\textbf{ID} & \textbf{Nome} & \textbf{Condição} & \textbf{Ação} & \textbf{Sever.} \\ \hline
\endfirsthead
\hline
\textbf{ID} & \textbf{Nome} & \textbf{Condição} & \textbf{Ação} & \textbf{Sever.} \\ \hline
\endhead

ROB\_001 & Mutex Global & \texttt{asset.has\_active\_workflow == true} & VETO & HIGH \\ \hline

ROB\_002 & Circuit Breaker Econômico & \texttt{token\_spend > daily\_limit} & DEGRADE & MEDIUM \\ \hline

ROB\_003 & Lockout/Tagout Virtual & \texttt{asset.loto\_flag == true} & VETO & BLOCKING \\ \hline

ROB\_004 & Dedupe de Eventos & \texttt{event.id} IN \texttt{processed\_cache} & LOG & LOW \\ \hline

ROB\_005 & Backpressure Management & \texttt{queue.depth > MAX\_QUEUE\_DEPTH} & DEGRADE & MEDIUM \\ \hline

ROB\_006 & Graceful Degradation & \texttt{llm.response\_time > 5000ms} & DEGRADE & MEDIUM \\ \hline

\caption{Guardrails de Robustez (ROB\_001 a ROB\_006)}
\end{longtable}
}

\subsection{Extensibilidade (Custom Guardrails)}

Os 38 guardrails listados acima constituem o \textbf{Catálogo Base} do sistema. No entanto, o motor de políticas foi desenhado para ser extensível, suportando regras customizadas em quatro níveis hierárquicos:

\begin{enumerate}
    \item \textbf{Guardrails por Tenant}: Regras de negócio específicas de um cliente (ex: "Cliente X proíbe manutenção aos domingos"). Sobrescrevem regras P2/P3, mas nunca P0/P1.
    \item \textbf{Guardrails por Ativo}: Regras vinculadas ao modelo do equipamento (ex: "Chiller Model Y tem limite térmico de 105°C ao invés de 110°C").
    \item \textbf{Guardrails Temporários}: Regras com TTL (Time-To-Live). Úteis para janelas de manutenção (ex: "Ignorar alertas de pressão por 2 horas durante teste").
    \item \textbf{Guardrails Emergentes}: Regras criadas dinamicamente pelo Viva Co-Pilot em resposta a incidentes novos (ex: "Enquanto não houver patch, bloquear comando Z").
\end{enumerate}

Essa arquitetura garante que o CAOS possa se adaptar a contratos específicos sem necessidade de alteração no código-fonte do núcleo.

% --------------------------------------------------------------------------
% RESUMO DO CATÁLOGO
% --------------------------------------------------------------------------
\subsubsection{Resumo do Catálogo}

\begin{table}[h!]
\centering
\begin{tabular}{|l|c|c|c|c|}
\hline
\textbf{Categoria} & \textbf{Qtd} & \textbf{BLOCKING} & \textbf{HIGH} & \textbf{MEDIUM/LOW} \\ \hline
Físicas (Safety) & 10 & 8 & 2 & 0 \\ \hline
Financeiras & 5 & 1 & 2 & 2 \\ \hline
Contratuais & 5 & 2 & 1 & 2 \\ \hline
Comunicação & 4 & 0 & 0 & 4 \\ \hline
Operacionais & 8 & 4 & 3 & 1 \\ \hline
Robustez & 6 & 1 & 1 & 4 \\ \hline
\textbf{Total} & \textbf{38} & \textbf{16} & \textbf{9} & \textbf{13} \\ \hline
\end{tabular}
\caption{Resumo do Catálogo de Guardrails}
\end{table}

\newpage

% ============================================================================
% CAPITULO 6: FLUXOS E CONTRATOS
% ============================================================================
\section{Fluxos Operacionais e Contratos}
\label{sec:fluxos}



\begin{figure}[h!]
    \centering
    \resizebox{\textwidth}{!}{%
    \begin{tikzpicture}[
        node distance=2cm,
        actor/.style={rectangle, draw=black, thick, fill=gray!10, minimum width=1.6cm, minimum height=0.9cm, align=center, font=\scriptsize},
        msg/.style={->, thick, >=stealth},
        msglabel/.style={font=\tiny, fill=white, inner sep=1pt},
        phase/.style={rectangle, draw=none, fill=gray!20, minimum width=2.2cm, minimum height=0.6cm, align=center, font=\scriptsize\bfseries, rounded corners=3pt}
    ]
    
    % Actors at y=0
    \node [actor] (sensor) at (0,0) {Sensor\\IoT};
    \node [actor] (pubsub) at (2.2,0) {Pub/Sub};
    \node [actor, fill=purple!15] (atlas) at (4.4,0) {ATLAS};
    \node [actor, fill=red!15] (sentinel) at (6.6,0) {SENTINEL};
    \node [actor, fill=cyan!15] (oracle) at (8.8,0) {ORACLE};
    \node [actor, fill=blue!20] (caos) at (11,0) {CAOS\\(Judge)};
    \node [actor, fill=green!15] (care) at (13.2,0) {CARE\\(HAL)};
    \node [actor] (device) at (15.4,0) {Device};
    
    % Lifelines
    \foreach \n in {sensor, pubsub, atlas, sentinel, oracle, caos, care, device} {
        \draw [dashed, gray] (\n.south) -- ++(0,-14);
    }
    
    % === FASE 1: INGESTÃO (Msgs 1-3) ===
    \node [phase, fill=purple!20, text width=2.5cm] at (-2.5,-1.75) {1. Ingestão\\(Msgs 1-3)};
    \draw [msg] (0,-1) -- node[msglabel, above] {1. Telemetry} (2.2,-1);
    \draw [msg] (2.2,-1.5) -- node[msglabel, above] {2. Raw Data} (4.4,-1.5);
    \draw [msg] (4.4,-2.2) -- ++(0.5,0) -- ++(0,-0.25) -- ++(-0.5,0) node[msglabel, right, xshift=0.5cm, yshift=0.1cm] {3. Update Twin};
    
    % === FASE 2: PERCEPÇÃO (Msgs 4-5) ===
    \node [phase, fill=red!20, text width=2.5cm] at (-2.5,-3.6) {2. Percepção\\(Msgs 4-5)};
    \draw [msg] (4.4,-3.2) -- node[msglabel, above] {4. State + Thresholds} (6.6,-3.2);
    \draw [msg] (6.6,-3.9) -- node[msglabel, above] {5. Anomaly Alert} (11,-3.9);
    
    % === FASE 3: PREDIÇÃO (Opcional - Msgs 6-7) ===
    \node [phase, fill=cyan!20, text width=2.5cm] at (-2.5,-5.0) {3. Predição\\(Opcional)};
    % 5.5. CAOS queries ATLAS for context (implicit, before reasoning)
    \draw [msg, dashed] (11,-4.3) -- node[msglabel, above] {Context Query} (4.4,-4.3);
    \draw [msg, dashed] (4.4,-4.6) -- node[msglabel, above] {RAG Memory} (11,-4.6);
    % 6. CAOS asks ORACLE for prediction
    \draw [msg, dashed] (11,-5.0) -- node[msglabel, above] {6. Simulate} (8.8,-5.0);
    \draw [msg, dashed] (8.8,-5.5) -- node[msglabel, above] {7. Forecast} (11,-5.5);
    
    % === FASE 4: JULGAMENTO HÍBRIDO (Msgs 8-11) ===
    \node [phase, fill=blue!20, text width=2.5cm] at (-2.5,-7.2) {4. Julgamento\\(Msgs 8-11)};
    % 4 sub-steps as self-loops on CAOS (labels in BLACK for readability)
    \draw [msg] (11,-6.2) -- ++(0.5,0) -- ++(0,-0.25) -- ++(-0.5,0) node[msglabel, right, xshift=0.5cm, yshift=0.1cm] {8. Analyze Impacts};
    \draw [msg] (11,-6.7) -- ++(0.5,0) -- ++(0,-0.25) -- ++(-0.5,0) node[msglabel, right, xshift=0.5cm, yshift=0.1cm] {9. Calculate S \& V};
    \draw [msg] (11,-7.2) -- ++(0.5,0) -- ++(0,-0.25) -- ++(-0.5,0) node[msglabel, right, xshift=0.5cm, yshift=0.1cm] {10. Check Guardrails};
    \draw [msg] (11,-7.7) -- ++(0.5,0) -- ++(0,-0.25) -- ++(-0.5,0) node[msglabel, right, xshift=0.5cm, yshift=0.1cm] {11. Classify Risk};
    
    % === FASE 5: EXECUÇÃO (Msgs 12-15) ===
    \node [phase, fill=green!20, text width=2.5cm] at (-2.5,-9.5) {5. Execução\\(Msgs 12-15)};
    \draw [msg] (11,-8.5) -- node[msglabel, above] {12. Execute Plan} (13.2,-8.5);
    \draw [msg] (13.2,-9.2) -- node[msglabel, above] {13. HAL Command} (15.4,-9.2);
    \draw [msg] (15.4,-9.9) -- node[msglabel, above] {14. ACK} (13.2,-9.9);
    \draw [msg] (13.2,-10.6) -- node[msglabel, above] {15. Result} (11,-10.6);
    
    % === FASE 6: COMUNICAÇÃO (Msgs 16-17) ===
    \node [phase, fill=orange!20, text width=2.5cm] at (-2.5,-12) {6. Comunicação\\(Msgs 16-17)};
    \draw [msg] (13.2,-11.3) -- ++(0.5,0) -- ++(0,-0.25) -- ++(-0.5,0) node[msglabel, right, xshift=0.5cm, yshift=0.1cm] {16. Notify Human};
    \draw [msg] (13.2,-12) -- ++(-11.2,0) node[msglabel, above, pos=0.5] {17. Status Update (Co-Pilot)};

    \end{tikzpicture}
    }
    \caption{Diagrama de Sequência: 6 Fases do Ciclo de Decisão Autônoma}
    \label{fig:sequence_diagram}
\end{figure}

\begin{table}[h!]
\centering
\small
\begin{tabular}{|c|c|l|l|l|}
\hline
\textbf{Fase} & \textbf{\#} & \textbf{De → Para} & \textbf{Mensagem} & \textbf{Latência} \\ \hline
\multirow{3}{*}{\textbf{1. Ingestão}} & 1 & Sensor → Pub/Sub & Telemetry JSON & <10ms \\ \cline{2-5}
& 2 & Pub/Sub → ATLAS & Raw Data & <20ms \\ \cline{2-5}
& 3 & ATLAS (interno) & Update Digital Twin & <50ms \\ \hline
\multirow{2}{*}{\textbf{2. Percepção}} & 4 & ATLAS → SENTINEL & State + Thresholds & <30ms \\ \cline{2-5}
& 5 & SENTINEL → CAOS & Anomaly Alert & <50ms \\ \hline
\multirow{4}{*}{\textbf{3. Predição (Opcional)}} & 5.1 & CAOS → ATLAS & Context Query & <30ms \\ \cline{2-5}
& 5.2 & ATLAS → CAOS & RAG Memory & <100ms \\ \cline{2-5}
& 6 & CAOS → ORACLE & Simulate & <50ms \\ \cline{2-5}
& 7 & ORACLE → CAOS & Forecast & <300ms \\ \hline
\multirow{4}{*}{\textbf{4. Julgamento}} & 8 & CAOS (interno) & Analyze Impacts & <50ms \\ \cline{2-5}
& 9 & CAOS (interno) & Calculate S \& V & <100ms \\ \cline{2-5}
& 10 & CAOS (interno) & Check Guardrails & <50ms \\ \cline{2-5}
& 11 & CAOS (interno) & Classify Risk & <20ms \\ \hline
\multirow{4}{*}{\textbf{5. Execução}} & 12 & CAOS → CARE & Execute Plan & <50ms \\ \cline{2-5}
& 13 & CARE → Device & HAL Command & <200ms \\ \cline{2-5}
& 14 & Device → CARE & ACK & <100ms \\ \cline{2-5}
& 15 & CARE → CAOS & Result & <50ms \\ \hline
\multirow{2}{*}{\textbf{6. Comunicação}} & 16 & CARE → Humano & Notify & <50ms \\ \cline{2-5}
& 17 & CARE → Co-Pilot & Status Update & <50ms \\ \hline
\multicolumn{2}{|c|}{\textbf{Total E2E}} & \multicolumn{2}{l|}{\textbf{Sensor → Notificação}} & \textbf{<1.5s} \\ \hline
\end{tabular}
\caption{Latências por Fase (17 Mensagens do Ciclo de Decisão)}
\end{table}

\subsection{Contratos de Interface}

O CAOS opera como um \textbf{event-driven agent}. Ele não expõe APIs REST, mas define contratos rígidos de entrada e saída.

\subsubsection{Input: Evento de Alerta (Sentinel)}

\begin{lstlisting}[language=Python, caption=Schema de Entrada - AlertEvent]
{
  "event_id": "evt_a1b2c3d4",
  "timestamp": "2026-01-28T10:00:00Z",
  "source": "sentinel",
  "asset_id": "CHILLER-04",
  "tenant_id": "tenant_abc",
  "severity": "CRITICAL",  // CRITICAL | HIGH | MEDIUM | LOW
  "metric": "temperature",
  "value": 95.2,
  "unit": "celsius",
  "threshold_violated": 90.0,
  "context": {
    "location": "Loja 042 - São Paulo",
    "model": "Carrier 30XA"
  }
}
\end{lstlisting}

\subsubsection{Output: Ação para Cloud Workflows (ActionSchema)}

\begin{lstlisting}[language=Python, caption=Schema de Saída - ActionSchema (destinado ao CARE)]
{
  "action_id": "act_x1y2z3",
  "workflow_name": "EmergencyShutdownWorkflow",
  "created_at": "2026-01-28T10:00:01Z",
  "verdict_score": -0.34,
  "decision": "BLOCKED",
  "asset_id": "CHILLER-04",
  "tenant_id": "tenant_abc",
  "command": {
    "type": "shutdown",
    "target": "compressor",
    "params": { "graceful": false }
  },
  "guardrails_passed": ["PHYS_001", "FIN_001"],
  "guardrails_violated": [],
  "requires_approval": false,
  "estimated_cost": 0.00,
  "justification": "Temperatura crítica. Risco de dano ao ativo."
}
\end{lstlisting}

\subsubsection{APIs Consumidas}

\begin{table}[h!]
\centering
\small
\begin{tabular}{|l|l|l|l|}
\hline
\textbf{Agente} & \textbf{Endpoint} & \textbf{Método} & \textbf{Descrição} \\ \hline
Atlas & \texttt{/api/v1/assets/\{id\}/twin} & GET & Digital Twin do ativo \\ \hline
Atlas & \texttt{/api/v1/assets/\{id\}/manuals} & GET & RAG de manuais técnicos \\ \hline
Sentinel & \texttt{/api/v1/alerts/active} & GET & Alertas ativos \\ \hline
Oracle & \texttt{/api/v1/predict/\{asset\_id\}} & POST & Predição de falha \\ \hline
\end{tabular}
\caption{APIs Consumidas pelo CAOS}
\end{table}

\subsection{Fluxo de Escalation (Viva Co-Pilot)}

Quando o Veredito requer aprovação humana (\texttt{REQUIRE\_APPROVAL}), o CAOS aciona o fluxo de escalation.

\subsubsection{Canais de Notificação}

\begin{table}[h!]
\centering
\small
\begin{tabular}{|l|l|l|l|}
\hline
\textbf{Canal} & \textbf{Severidade} & \textbf{Timeout} & \textbf{Auto-Approve} \\ \hline
WhatsApp & CRITICAL, HIGH & 30min & Não \\ \hline
Dashboard & Todos & 4h & Sim (após timeout) \\ \hline
Email & MEDIUM, LOW & 24h & Sim (após timeout) \\ \hline
SMS & CRITICAL only & 15min & Não \\ \hline
\end{tabular}
\caption{Matriz de Canais de Escalation}
\end{table}

\subsubsection{Estados de Aprovação}

\begin{enumerate}
    \item \textbf{PENDING}: Aguardando resposta humana.
    \item \textbf{APPROVED}: Humano aprovou a ação.
    \item \textbf{REJECTED}: Humano rejeitou a ação.
    \item \textbf{TIMEOUT}: Tempo expirou. Aplica política padrão.
    \item \textbf{AUTO\_APPROVED}: Timeout + política permite auto-approve.
\end{enumerate}

\newpage

\subsection{Políticas JSON Completas}

Abaixo, exemplos de políticas para cada categoria de guardrail.

\subsubsection{physical.json}
\begin{lstlisting}[language=Python, caption=Política Física, basicstyle=\tiny\ttfamily]
{
  "policy_id": "PHYS",
  "version": "1.0",
  "rules": [
    { "id": "PHYS_001", "condition": "temp_c > 110", "action": "VETO", "severity": "BLOCKING" },
    { "id": "PHYS_004", "condition": "fire_detected == true", "action": "VETO", "severity": "BLOCKING" },
    { "id": "PHYS_008", "condition": "under_maintenance == true", "action": "VETO", "severity": "BLOCKING" }
  ]
}
\end{lstlisting}

\subsubsection{financial.json}
\begin{lstlisting}[language=Python, caption=Política Financeira, basicstyle=\tiny\ttfamily]
{
  "policy_id": "FIN",
  "version": "1.0",
  "rules": [
    { "id": "FIN_001", "condition": "action.cost > 500", "action": "REQUIRE_APPROVAL", "role": "MANAGER" },
    { "id": "FIN_003", "condition": "monthly_spend > budget", "action": "VETO", "severity": "BLOCKING" }
  ]
}
\end{lstlisting}

\subsubsection{contractual.json}
\begin{lstlisting}[language=Python, caption=Política Contratual, basicstyle=\tiny\ttfamily]
{
  "policy_id": "CONT",
  "version": "1.0",
  "rules": [
    { "id": "CONT_001", "condition": "action NOT IN contract.allowed", "action": "VETO" },
    { "id": "CONT_004", "condition": "tier == 'MONITORING_ONLY'", "action": "VETO" }
  ]
}
\end{lstlisting}

\subsubsection{communication.json}
\begin{lstlisting}[language=Python, caption=Política de Comunicação, basicstyle=\tiny\ttfamily]
{
  "policy_id": "COMM",
  "version": "1.0",
  "rules": [
    { "id": "COMM_001", "condition": "sentiment NOT IN ['professional']", "action": "REQUIRE_APPROVAL" },
    { "id": "COMM_002", "condition": "same_alert_count > 5 PER MINUTE", "action": "LOG" }
  ]
}
\end{lstlisting}

\subsubsection{secops.json}
\begin{lstlisting}[language=Python, caption=Política SecOps, basicstyle=\tiny\ttfamily]
{
  "policy_id": "SO",
  "version": "1.0",
  "rules": [
    { "id": "SO_001", "condition": "decision_changed IN < 30s", "action": "VETO" },
    { "id": "SO_002", "condition": "payload.contains(INJECTION)", "action": "VETO", "severity": "BLOCKING" }
  ]
}
\end{lstlisting}

\subsubsection{robustness.json}
\begin{lstlisting}[language=Python, caption=Política de Robustez, basicstyle=\tiny\ttfamily]
{
  "policy_id": "ROB",
  "version": "1.0",
  "rules": [
    { "id": "ROB_001", "condition": "asset.has_active_workflow", "action": "VETO" },
    { "id": "ROB_003", "condition": "loto_flag == true", "action": "VETO", "severity": "BLOCKING" }
  ]
}
\end{lstlisting}

\newpage

\subsection{Métricas e Observabilidade}

\subsubsection{Dashboards (Google Cloud Monitoring)}

\begin{table}[h!]
\centering
\small
\begin{tabular}{|l|l|l|}
\hline
\textbf{Dashboard} & \textbf{Métricas} & \textbf{Alerta} \\ \hline
CAOS Latency & P50, P95, P99 de decisão & P99 > 2s \\ \hline
CAOS Throughput & Decisões/min por tenant & < 10/min warn \\ \hline
Guardrail Violations & Contagem por categoria & > 50/hora \\ \hline
Workflows Health & Workflows pending/failed & Failed > 5 \\ \hline
LLM Token Usage & Tokens/dia, custo USD & > \$100/dia \\ \hline
\end{tabular}
\caption{Dashboards de Observabilidade}
\end{table}

\subsubsection{Latency Targets (SLOs)}

\begin{itemize}
    \item \textbf{P50}: < 500ms (decisão típica)
    \item \textbf{P95}: < 1.5s (decisão com RAG)
    \item \textbf{P99}: < 3s (edge cases com múltiplos agentes)
    \item \textbf{Availability}: 99.9\% uptime mensal
\end{itemize}

\subsection{Rollback e Recovery}

\subsubsection{Cenários de Falha e Recuperação}

\begin{table}[h!]
\centering
\small
\begin{tabular}{|l|l|l|}
\hline
\textbf{Falha} & \textbf{Comportamento} & \textbf{Recovery} \\ \hline
Cloud Workflows down & Workflows pausados & Auto-resume ao subir \\ \hline
Firestore down & State read fails & Fallback para cache Redis \\ \hline
Vertex AI quota & LLM calls fail & Fallback determinístico \\ \hline
Pub/Sub timeout & Events queued & Retry com backoff \\ \hline
Execution error & Step fails & Workflow retry (3x) \\ \hline
\end{tabular}
\caption{Matriz de Falha e Recovery}
\end{table}

\subsubsection{Idempotência}

Todas as ações do CAOS são \textbf{idempotentes}. Se um comando \texttt{shutdown} for enviado 2x, o Worker verifica o estado atual antes de executar. Se o device já está desligado, retorna \texttt{NOOP}.

\subsection{Testes de Integração}

\subsubsection{Framework de Testes}
\begin{itemize}
    \item \textbf{Unit Tests}: pytest + mocks para LangGraph nodes
    \item \textbf{Integration Tests}: Cloud Workflows Emulator + fixtures
    \item \textbf{E2E Tests}: Pub/Sub emulator + real LLM calls (staging)
\end{itemize}

\subsubsection{Cenários de Teste Automatizados}

\begin{table}[h!]
\centering
\small
\begin{tabular}{|l|l|l|}
\hline
\textbf{Cenário} & \textbf{Input} & \textbf{Expected Output} \\ \hline
Fire Detection & severity=CRITICAL, fire=true & V < 0, VETO \\ \hline
Normal Operation & severity=LOW, all OK & V > 0.5, MONITOR \\ \hline
Budget Exceeded & cost > budget & REQUIRE\_APPROVAL \\ \hline
LOTO Active & under\_maintenance=true & VETO \\ \hline
Prompt Injection & payload="ignore all" & VETO, BLOCKED \\ \hline
\end{tabular}
\caption{Cenários de Teste Automatizados}
\end{table}

\subsection{Rate Limits e Graceful Degradation}

\subsubsection{Limites do Vertex AI (Gemini)}

\begin{table}[h!]
\centering
\small
\begin{tabular}{|l|l|l|}
\hline
\textbf{Recurso} & \textbf{Limite} & \textbf{Ação se Excedido} \\ \hline
Requests/min & 60 RPM (default) & Queue + backoff \\ \hline
Tokens/day & 2M tokens & Switch to cache \\ \hline
Cost/day & \$50 (configurável) & Fallback determinístico \\ \hline
\end{tabular}
\caption{Limites de API e Ações}
\end{table}

\subsubsection{Fallback Determinístico}

Quando a API do LLM está indisponível, o CAOS entra em \textbf{Modo Degradado}:
\begin{enumerate}
    \item \textbf{Alertas CRITICAL}: Executa regras hardcoded (shutdown automático)
    \item \textbf{Alertas HIGH}: Envia para Co-Pilot sem análise LLM
    \item \textbf{Alertas MEDIUM/LOW}: Loga e ignora até recovery
\end{enumerate}

\newpage


% ============================================================================
% CAPITULO 7: SINAPSES (PROTOCOLOS)
% ============================================================================
\section{Integração com Agentes}
\label{sec:sinapses}

\subsection{Topologia (Google Pub/Sub)}
\begin{table}[h!]
\centering
\begin{tabular}{|l|l|l|}
\hline
\textbf{Canal} & \textbf{De -> Para} & \textbf{Conteúdo} \\ \hline
\texttt{sentinel.alerts} & Sentinel -> CAOS & Trigger JSON \\ \hline
\texttt{oracle.predictions} & Oracle -> CAOS & Probabilidades \\ \hline
\texttt{caos.actions} & CAOS -> Care & Command Object \\ \hline
\end{tabular}
\end{table}

\subsection{Payloads Estritos}
\begin{lstlisting}[language=Python, caption=Exemplo de Trigger]
{
  "source": "sentinel", 
  "payload": { "asset_id": "PUMP-04", "metric": "vibration", "val": 12.5 }
}
\end{lstlisting}

\newpage



% ============================================================================
% CAPITULO 8: INTERFACES (API & TOOLS)
% ============================================================================
\section{Ferramentas (Tools)}
\label{sec:tools}

O Cortex não apenas "fala", ele "faz". As seguintes funções estão disponíveis para o LLM via \textit{Function Calling}:

\subsection{Catálogo de Tools}

\begin{table}[h!]
\centering
\small
\begin{tabular}{|l|l|p{6cm}|}
\hline
\textbf{Tool} & \textbf{Parâmetros} & \textbf{Descrição} \\ \hline
\texttt{mqtt\_publish} & topic, payload, qos & Envia comando direto a device IoT via broker MQTT \\ \hline
\texttt{open\_ticket} & severity, queue, title, description & Abre chamado no Jira/ServiceNow \\ \hline
\texttt{send\_whatsapp} & user\_id, template\_id, params & Notificação urgente via Twilio \\ \hline
\texttt{sql\_query} & query (read-only) & Consulta analítica no TimescaleDB \\ \hline
\texttt{get\_asset\_state} & asset\_id & Recupera Digital Twin do Atlas \\ \hline
\texttt{schedule\_maintenance} & asset\_id, datetime, type & Agenda manutenção preventiva \\ \hline
\end{tabular}
\caption{Catálogo de Ferramentas do CAOS}
\end{table}

\subsection{Exemplo de Uso (Function Calling)}

\begin{lstlisting}[language=Python, caption=Exemplo de Tool Call pelo LLM, basicstyle=\tiny\ttfamily]
{
  "tool": "mqtt_publish",
  "parameters": {
    "topic": "devices/chiller-04/commands",
    "payload": {"action": "shutdown", "graceful": true},
    "qos": 1
  },
  "reasoning": "Temperatura critica detectada. Shutdown preventivo."
}
\end{lstlisting}

\newpage

\section{API REST}
\label{sec:api}

O CAOS expõe endpoints via FastAPI para integração com sistemas legados (ERPs) ou Dashboards.

\begin{table}[h!]
\centering
\begin{tabular}{|l|l|l|}
\hline
\textbf{Método} & \textbf{Endpoint} & \textbf{Função} \\ \hline
POST & \texttt{/v1/event} & Ingestão Manual de Eventos \\ \hline
GET & \texttt{/v1/trace/\{id\}} & Consulta de Cadeia de Pensamento \\ \hline
POST & \texttt{/v1/feedback} & Reinforcement Learning Human Feedback \\ \hline
\end{tabular}
\end{table}

\subsection{Exemplo de Requisicão (Ingestão)}
\begin{lstlisting}[language=bash]
curl -X POST https://api.vivaiot.com/v1/event \
  -H "Authorization: Bearer <token>" \
  -d '{
    "source": "maintenance_app",
    "type": "alert_ack",
    "payload": {
       "alert_id": "ALT-99",
       "user": "technician_01",
       "action": "investigating"
    }
  }'
\end{lstlisting}

\subsection{Exemplo de Resposta (Trace)}
\begin{lstlisting}[language=Python]
{
  "trace_id": "abc-123",
  "status": "processing",
  "steps": [
    {"agent": "Sentinel", "output": "Anomaly detected"},
    {"agent": "Atlas", "output": "Context found: Maintenance Window"},
    {"agent": "Oracle", "output": "Decision: Suppress Alert"}
  ]
}
\end{lstlisting}

\newpage

% ============================================================================
% CAPITULO 9: SEGURANÇA (ENTERPRISE)
% ============================================================================
\section{Segurança e Conformidade}
\label{sec:security}

\subsection{Proteção na Borda (mTLS)}
O Gateway Edge e a Nuvem CAOS utilizam \textbf{Mutual TLS}. Ambos os lados apresentam certificados assinados pela CA interna. Tokens de execução têm TTL de 5 minutos.

\subsection{Provas Jurídicas (WORM)}
Os logs de veredito são armazenados em buckets S3 com \textbf{Object Lock} (WORM) por 5 anos. Garantia de imutabilidade para auditoria.

\newpage

% ============================================================================
% CAPITULO 10: RESILIENCIA (DR)
% ============================================================================
\section{Resiliência e Recuperação}
\label{sec:dr}

\subsection{Protocolo Brain Death}
Se o LLM falhar (timeout > 5s):
\begin{enumerate}
    \item \textbf{Fallback}: CAOS entra em "Modo Reptiliano".
    \item \textbf{Lógica}: Ignora Oracle. Usa apenas Sentinel + Guardrails.
    \item \textbf{Restrição}: Apenas ações de Segurança (P1) são permitidas.
\end{enumerate}

\subsection{Circuit Breaker}
Se um ativo falhar o reset 3x, o circuito abre por 4 horas.

\newpage

% ============================================================================
% CAPITULO 11: ESTRATEGIA DE TESTES (QA)
% ============================================================================
\section{Estratégia de Qualidade}
\label{sec:qa}

\subsection{Pirâmide de Testes}
\begin{enumerate}
    \item \textbf{Unitários (80\%)}: Foco nos Guardrails (Input fixo -> Block/Allow).
    \item \textbf{Integração (15\%)}: Teste de conversores HAL (Modbus -> JSON).
    \item \textbf{End-to-End (5\%)}: Simulação completa via "Digital Twin" (Mock do Sentinel).
\end{enumerate}

\newpage

% ============================================================================
% CAPITULO 12: OPERAÇÃO E OBSERVABILIDADE OPS
% ============================================================================
\section{Observabilidade}
\label{sec:ops}

Como garantir que o sistema está saudável sem olhar os logs manualmente?

\subsection{Tracing Cognitivo (LangSmith)}
Diferente de sistemas comuns, o CAOS precisa de debugging *semântico*. Utilizamos o \textbf{LangSmith} para visualizar a Cadeia de Pensamento.

\begin{figure}[h!]
    \centering
    \begin{tcolorbox}[title=Exemplo de Trace Log, colback=black!5!white]
    \small
    \textbf{[Trace ID: 9f8a...]} \\
    \texttt{> User Input: "Vibração alta na Bomba 1"} \\
    \texttt{  |-- Agent: Sentinel (Output: Severity HIGH)} \\
    \texttt{  |-- Agent: Atlas (Tool: vector\_search "manual bomba 1")} \\
    \texttt{  |   |-- Chunk 1: "Vibração acima de 80\% requer parada."} \\
    \texttt{  |-- Agent: Oracle (Thinking...)} \\
    \texttt{  |-- Agent: Oracle (Decision: STOP\_MACHINE)} \\
    \texttt{  |-- Guardrail: "Stop Machine" is SAFE? -> YES} \\
    \texttt{< System Action: mqtt\_publish("STOP")}
    \end{tcolorbox}
    \caption{Visualização de um Trace Cognitivo}
\end{figure}

\subsection{Infraestrutura-como-Código (Terraform)}
Todo o ambiente Google Cloud (Vertex AI, Cloud Run, PubSub, Postgres) é provisionado via Terraform.

\begin{lstlisting}[language=bash, caption=Exemplo de Recurso Terraform]
resource "google_cloud_run_service" "caos_agent" {
  name     = "caos-agent-v1"
  location = "us-central1"

  template {
    spec {
      containers {
        image = "gcr.io/vivaiot/caos-agent:latest"
        env {
          name = "PUBSUB_PROJECT"
          value = "vivaiot-prod-001"
        }
      }
    }
  }
}
\end{lstlisting}

\newpage
\section{Deploy}
O CAOS roda em containers \textit{Serverless}:
\begin{itemize}
    \item \textbf{Escala Zero}: Se não houver alertas, custo é zero (exceto DB).
    \item \textbf{Isolamento}: Cada Tenant (Cliente) roda em uma instância isolada.
\end{itemize}

\subsection{Dockerfile de Produção}
\begin{lstlisting}[language=bash]
FROM python:3.11-slim

WORKDIR /app
COPY requirements.txt .
RUN pip install --no-cache-dir -r requirements.txt

COPY src/ .
CMD ["uvicorn", "caos.api.main:app", "--host", "0.0.0.0", "--port", "8080"]
\end{lstlisting}


\newpage

% ============================================================================
% CAPITULO 12: MATRIZ DE CENARIOS
% ============================================================================
\section{Exemplos de Fluxo}
\label{sec:scenarios}

\subsection{Indústria (Compressor Superaquecido)}

\begin{enumerate}
    \item \textbf{Evento}: Sensor IoT reporta temperatura de 105°C (Limite 100°C).
    \item \textbf{Sentinel}: Classifica como "Crítico" (Score 0.98).
    \item \textbf{Atlas}: Busca manual. "Reset permitido se < 110°C". Busca agenda: "Produção ativa".
    \item \textbf{Oracle (se acionado)}: Decide "Tentar Reset Remoto antes de parar linha" para otimização de custo.
    \item \textbf{Guardrail}: Checa "Reset Remoto". Seguro? Sim.
    \item \textbf{Action}: \texttt{mqtt\_publish("cmd/reset", \{"force": true\})}
    \item \textbf{Resultado}: Temperatura cai para 95°C. Ticket de "Atenção" aberto.
\end{enumerate}

\subsection{Varejo (Smart Fridge)}
\textbf{Cenário}: Porta aberta por 10min.
\begin{itemize}
    \item \textbf{Decisão}: Não desligar (prejuízo de estoque).
    \item \textbf{Ação}: Enviar WhatsApp para Gerente da Loja.
    \item \textbf{Output}: \texttt{send\_whatsapp("Marta", "Geladeira 4 aberta!")}
\end{itemize}

\subsection{Segurança (Fraude)}
\textbf{Cenário}: Tentativa de acesso com cartão clonado.
\begin{itemize}
    \item \textbf{Sentinel}: Padrão de anomalia geo-espacial.
    \item \textbf{Guardrail}: Bloqueio Imediato (P1).
    \item \textbf{Ação}: Bloquear porta + Acionar Câmera.
\end{itemize}

\newpage

% ============================================================================
% CAPITULO 13: ESPECIFICAÇÕES DE IMPLEMENTAÇÃO (CÓDIGO)
% ============================================================================
\section{Especificações de Implementação}
\label{sec:specs}

Esta seção define a estrutura física do projeto Python e os padrões de desenvolvimento.

\subsection{Estrutura de Diretórios}
\begin{lstlisting}[language=bash]
src/
|-- caos/
|   |-- core/           # O "Cérebro"
|   |   |-- brain.py    # LangGraph StateGraph
|   |   |-- judge.py    # Lógica de Veredicto (Fórmula)
|   |-- integration/    # Clientes de Serviços Externos
|   |   |-- atlas_client.py     # Consulta Digital Twin
|   |   |-- workflows_client.py # Disparo para o Agente CARE
|   |-- safety/            # O "Sistema Imunológico"
|   |   |-- guardrail.py   # Validador de Policies
|-- actions/               # Camada CARE (Cloud Run/Functions)
|   |-- action_api.py      # Executor HAL
|   |-- steps/
|   |   |-- iot_actions.py  # (Ex: Mqtt Publish)
|   |   |-- notifications.py
|   |-- workflows/
|   |   |-- remediation.yaml # Workflow YAML (CARE engine)
|-- tests/              # Pirâmide de Testes (Unit/Int/E2E)
|-- main.py             # Entrypoint (FastAPI)
\end{lstlisting}

\subsection{Interfaces Centrais (Abstrações)}

\textbf{1. O Protocolo HAL (IActuator)}
Todo driver deve implementar esta interface para garantir desacoplamento.
\begin{lstlisting}[language=Python]
class IActuator(ABC):
    @abstractmethod
    async def execute(self, command: Command) -> Result:
        """Envia sinal físico ao hardware"""
        pass
        
    @abstractmethod
    async def health_check(self) -> bool:
        """Heartbeat do driver"""
        pass
\end{lstlisting}

\textbf{2. O Contrato Cognitivo (JudgeState)}
O estado que flui pelo LangGraph e persiste no checkpointer.
\begin{lstlisting}[language=Python]
class JudgeState(TypedDict):
    trigger: TriggerPayload      # Gatilho (Sentinel/Oracle)
    atlas_context: AssetContext  # Contexto (Digital Twin/Docs)
    risk_dimensions: Dict[str, float] # {R_F, R_Fin, R_C, R_K}
    severity_score: float        # Escore S (0.0 a 1.0)
    verdict_score: float         # Escore V (-1.0 a 1.0)
    reasoning_trace: List[str]   # Chain-of-Thought
    proposed_action: Action      # Acao sugerida pelo LLM
    guardrail_violations: List[str] # Lista de bloqueios
\end{lstlisting}

\newpage

% ============================================================================
% CAPITULO 14: LIMITAÇÕES CONHECIDAS
% ============================================================================
\section{Limitações Conhecidas}
\label{sec:limitations}

Esta seção documenta as limitações atuais do sistema CAOS, garantindo transparência sobre cenários não cobertos ou parcialmente suportados.

\subsection{Limitações Técnicas}

\begin{enumerate}
    \item \textbf{Dependência de LLM}: O sistema depende da disponibilidade do Vertex AI (Gemini). Em caso de indisponibilidade prolongada (>5min), apenas ações P0/P1 são processadas via fallback determinístico.
    
    \item \textbf{Latência de Predição}: O Oracle adiciona 200-500ms ao fluxo. Para alertas CRITICAL (<1s SLA), a predição pode ser bypassada.
    
    \item \textbf{Cold Start}: Containers Cloud Run levam 2-5s para inicializar. Alertas durante cold start ficam enfileirados no Pub/Sub.
    
    \item \textbf{Versionamento de Modelo}: Mudanças no modelo LLM (ex: Gemini 1.5 → 2.0) podem alterar comportamento de decisão. Requer re-calibração de pesos.
    
    \item \textbf{Offline Edge}: Dispositivos IoT sem conectividade não recebem comandos até reconexão. O CAOS não possui cache local.
\end{enumerate}

\subsection{Limitações de Escopo}

\begin{enumerate}
    \item \textbf{Multi-Tenant Isolation}: Cada tenant roda em instância isolada, mas não há suporte a hierarquias complexas (ex: holding → subsidiárias).
    
    \item \textbf{Idioma}: O sistema é otimizado para Português (BR). Suporte a outros idiomas requer ajuste de prompts e guardrails de comunicação.
    
    \item \textbf{Protocolos Industriais}: Suporte nativo a MQTT e Modbus. Outros protocolos (BACnet, OPC-UA) requerem adaptadores customizados no HAL.
    
    \item \textbf{Explainability}: A justificativa do LLM é textual. Não há visualização gráfica do raciocínio (ex: árvore de decisão).
\end{enumerate}

\subsection{Edge Cases Não Cobertos}

\begin{table}[h!]
\centering
\small
\begin{tabular}{|l|l|l|}
\hline
\textbf{Cenário} & \textbf{Comportamento Atual} & \textbf{Mitigação Futura} \\ \hline
Conflito entre guardrails & Prioridade por severidade & Resolver via mediador \\ \hline
Ataque adversarial ao LLM & Bloqueio por prompt injection & Red team contínuo \\ \hline
Falha simultânea de 3+ serviços & Sistema para & Circuit breaker global \\ \hline
Rollback de ação física & Não suportado & Implementar compensação \\ \hline
\end{tabular}
\caption{Edge Cases e Mitigações Futuras}
\end{table}

\newpage

% ============================================================================
% CAPITULO 15: REFERENCIAS
% ============================================================================
\section{Referências}
\begin{itemize}
    \item LangGraph Documentation - \url{https://langchain-ai.github.io/langgraph/}
    \item Google Cloud Pub/Sub Documentation - \url{https://cloud.google.com/pubsub/docs}
    \item Google Cloud Workflows Documentation - \url{https://cloud.google.com/workflows/docs}
    \item Vertex AI Gemini API - \url{https://cloud.google.com/vertex-ai/docs/generative-ai/}
\end{itemize}

\end{document}

